\documentclass[a4paper,11pt]{amsart}

% =========================================================================
%  Informe de laboratorio — Práctica 3: Péndulos acoplados
%  Estilo: el del informe de Steiner (artículo científico minimalista,
%  versalitas, enunciados numerados por sección, cajas de referencia y de
%  método), con dos variaciones: tipografía Libertinus, más propia de los
%  artículos científicos, y una paleta de azul marino, azul petróleo,
%  turquesa y celeste.
%
%  Partes pendientes (marcadas en gris con \pend y \pendiente):
%  incertidumbres, ajustes lineales de la cuestión 2, comparación de la
%  cuestión 3 y conclusiones.
% =========================================================================

% --- IDIOMA Y CODIFICACIÓN ------------------------------------------------
\usepackage[utf8]{inputenc}
\usepackage[T1]{fontenc}
\usepackage[spanish,es-tabla,es-nodecimaldot]{babel}

% --- TIPOGRAFÍA (Libertinus: texto y matemáticas) ---------------------------
\usepackage{libertinus}
\linespread{1.06}
\usepackage[final]{microtype}

% --- MATEMÁTICAS Y UNIDADES ------------------------------------------------
\usepackage{amsmath, amssymb, mathtools}
\usepackage{bm}
\usepackage{siunitx}
\sisetup{
  output-decimal-marker = {,},
  group-minimum-digits  = 4,
  group-digits          = integer,
  group-separator       = {\,},
  per-mode              = symbol,
  separate-uncertainty  = true,
  range-phrase          = {\text{ a }},
  detect-all            = false
}

% --- PÁGINA, COLORES Y GRÁFICOS --------------------------------------------
\usepackage[a4paper,margin=2.6cm,footskip=1.45cm]{geometry}
\usepackage{graphicx}
\usepackage{xcolor}
\definecolor{azulprof}{HTML}{0B2545}    % azul marino
\definecolor{azulacero}{HTML}{1D5C7A}   % azul petróleo
\definecolor{turquesa}{HTML}{1E9E9B}    % turquesa
\definecolor{azulclaro}{HTML}{A8D5E2}   % celeste
\definecolor{grispizarra}{HTML}{4B5B6B} % gris pizarra azulado
\definecolor{grisclaro}{HTML}{E2EAEF}   % gris perla
\definecolor{azulfondo}{HTML}{F1F8FA}   % fondo tenue de referencias y resultados
\colorlet{acento}{turquesa!85!azulprof} % curvas teóricas y comparaciones
\usepackage{tikz}
\usetikzlibrary{babel, arrows.meta, calc, positioning, decorations.pathmorphing,
                decorations.pathreplacing, decorations.markings, patterns, shadings}
\usepackage{pgfplots}
\usepackage{pgfplotstable}
\usepgfplotslibrary{groupplots, fillbetween}
\pgfplotsset{compat=1.17}
\usepackage{booktabs}
\usepackage{multirow}
\usepackage[most]{tcolorbox}
\usepackage{array}
\usepackage{enumitem}
\setlist[itemize]{leftmargin=1.6em, itemsep=2pt, topsep=4pt}
\setlist[enumerate]{leftmargin=1.8em, itemsep=2pt, topsep=4pt}

% --- ENLACES Y REFERENCIAS CRUZADAS ----------------------------------------
\usepackage[colorlinks=true, linkcolor=azulacero, citecolor=azulacero,
            urlcolor=azulacero, plainpages=false, pdfusetitle=false,
            pdftitle={Práctica 3: Péndulos acoplados},
            pdfauthor={David Gómez García; Daniel Martín Guerrero}]{hyperref}
\usepackage[spanish,nameinlink]{cleveref}

% --- BIBLIOGRAFÍA -----------------------------------------------------------
\usepackage{csquotes}
\usepackage[backend=biber, style=numeric, sorting=none, giveninits=true,
            maxbibnames=6, minbibnames=3, doi=true, url=true,
            isbn=false]{biblatex}
\addbibresource{referencias.bib}
\DeclareFieldFormat[article,techreport,unpublished,misc]{title}{#1\isdot}
\DeclareFieldFormat{url}{\url{#1}}
\renewcommand*{\bibfont}{\small}
\setlength{\bibitemsep}{3pt}

% --- PIEZAS DE DIBUJO DE LOS PÉNDULOS ----------------------------------------
%%%%%%%%%%%%%%%%%%%%%%%%%%%%%%%%%%%%%%%%%%%%%%%%%%%%%%%%%%%%%%%%%%%%%%%%
% Piezas de dibujo de los péndulos acoplados (TikZ)
%
% Vista frontal y minimalista del montaje de la práctica: dos cabezales
% triangulares (como los soportes PHYWE de las fotos), dos varillas de
% longitud L con su masa lenticular y un muelle helicoidal sujeto a las dos
% varillas a una distancia \ell de los ejes de giro.
% Se usan en la portada y en la figura del montaje experimental.
% Requiere los colores azulprof, azulacero, turquesa, azulclaro,
% grispizarra y grisclaro (definidos en main.tex).
%%%%%%%%%%%%%%%%%%%%%%%%%%%%%%%%%%%%%%%%%%%%%%%%%%%%%%%%%%%%%%%%%%%%%%%%

% Punto a distancia #2 del pivote #1 a lo largo de una varilla girada #3 grados
\newcommand{\sobrevarilla}[3]{($(#1)+({#2*sin(#3)},{-#2*cos(#3)})$)}

% Cabezal triangular invertido con el eje de giro en el vértice inferior
% \cabezal{coordenada}{escala}
\newcommand{\cabezal}[2]{%
  \fill[azulprof] ($(#1)+(-0.62*#2,0.78*#2)$) -- ($(#1)+(0.62*#2,0.78*#2)$)
    -- ($(#1)+(0,-0.06*#2)$) -- cycle;
  \fill[azulacero!70!azulprof] ($(#1)+(-0.62*#2,0.78*#2)$) -- ($(#1)+(0.62*#2,0.78*#2)$)
    -- ($(#1)+(0.62*#2,0.86*#2)$) -- ($(#1)+(-0.62*#2,0.86*#2)$) -- cycle;
  \draw[white, line width=0.5pt, fill=azulprof] (#1) circle (0.085*#2);
  \fill[white] ($(#1)+(0,0.46*#2)$) circle (0.05*#2);}

% Muelle helicoidal entre dos puntos, con sus pinzas
% \muelleacople{A}{B}{escala}
\newcommand{\muelleacople}[3]{%
  \draw[azulacero, line width=0.55pt,
        decorate, decoration={coil, aspect=0.55, amplitude=3.1pt*#3,
                              segment length=2.9pt*#3, pre length=3pt*#3, post length=3pt*#3}]
    (#1) -- (#2);
  \fill[grispizarra] ($(#1)+(-0.07*#3,-0.07*#3)$) rectangle ($(#1)+(0.07*#3,0.07*#3)$);
  \fill[grispizarra] ($(#2)+(-0.07*#3,-0.07*#3)$) rectangle ($(#2)+(0.07*#3,0.07*#3)$);}

% Masa lenticular: disco con un brillo discreto
% \masalenticular{coordenada}{radio}
\newcommand{\masalenticular}[2]{%
  \shade[ball color=turquesa!85!white, opacity=1] (#1) circle (#2);
  \draw[azulprof, line width=0.6pt] (#1) circle (#2);
  \fill[white, opacity=0.35] ($(#1)+(-0.32*#2,0.34*#2)$) circle (0.22*#2);}

% El par de péndulos acoplados
% \parpendulos{phi1}{phi2}{ell}{L}{d}{escala}
%   phi1, phi2 : ángulos (grados) respecto a la vertical, positivos hacia la derecha
%   ell, L     : distancia del eje al muelle y longitud de la varilla
%   d          : separación entre los ejes de giro
%   escala     : tamaño de cabezales, masas y muelle
% Define las coordenadas P1, P2 (ejes), S1, S2 (muelle) y B1, B2 (masas).
\newcommand{\parpendulos}[6]{%
  \coordinate (P1) at ({-#5/2},0);
  \coordinate (P2) at ({#5/2},0);
  \coordinate (S1) at ($(P1)+({#3*sin(#1)},{-#3*cos(#1)})$);
  \coordinate (S2) at ($(P2)+({#3*sin(#2)},{-#3*cos(#2)})$);
  \coordinate (B1) at ($(P1)+({#4*sin(#1)},{-#4*cos(#1)})$);
  \coordinate (B2) at ($(P2)+({#4*sin(#2)},{-#4*cos(#2)})$);
  % verticales de equilibrio
  \draw[grispizarra!45, line width=0.45pt, dash pattern=on 3pt off 2.2pt]
    (P1) -- ++(0,{-#4-0.55*#6});
  \draw[grispizarra!45, line width=0.45pt, dash pattern=on 3pt off 2.2pt]
    (P2) -- ++(0,{-#4-0.55*#6});
  % varillas
  \draw[azulprof, line width=1.05pt*#6, line cap=round] (P1) -- (B1);
  \draw[azulprof, line width=1.05pt*#6, line cap=round] (P2) -- (B2);
  % muelle de acoplamiento
  \muelleacople{S1}{S2}{#6}
  % masas lenticulares y cabezales
  \masalenticular{B1}{0.36*#6}
  \masalenticular{B2}{0.36*#6}
  \cabezal{P1}{#6}
  \cabezal{P2}{#6}}

% Arco del ángulo girado por una varilla, con su rótulo. El rótulo se coloca
% al otro lado de la vertical de equilibrio, para que no pise la varilla.
% \angulo{pivote}{phi}{radio}{rótulo}
\newcommand{\angulo}[4]{%
  \draw[turquesa!80!azulprof, line width=0.6pt, -{Stealth[length=3.4pt]}]
    ($(#1)+(0,-#3)$) arc[start angle=-90, end angle={-90+#2}, radius=#3];
  \ifdim #2pt>0pt
    \node[text=azulprof, font=\small, anchor=east, inner sep=1.5pt] at ($(#1)+(-0.06,{-#3})$) {#4};
  \else
    \node[text=azulprof, font=\small, anchor=west, inner sep=1.5pt] at ($(#1)+(0.06,{-#3})$) {#4};
  \fi}

% Posición «fantasma» de un péndulo (el extremo opuesto de su oscilación)
% \fantasma{pivote}{phi}{L}{escala}
\newcommand{\fantasma}[4]{%
  \draw[azulclaro!90!grispizarra, line width=0.8pt*#4, line cap=round, opacity=0.8]
    (#1) -- ($(#1)+({#3*sin(#2)},{-#3*cos(#2)})$);
  \draw[azulclaro!90!grispizarra, line width=0.5pt, fill=azulfondo, opacity=0.9]
    ($(#1)+({#3*sin(#2)},{-#3*cos(#2)})$) circle (0.36*#4);}

% Flecha de movimiento a lo largo de la trayectoria de la masa
% \movimiento{pivote}{desde}{hasta}{radio}  (ángulos respecto a la vertical)
\newcommand{\movimiento}[4]{%
  \draw[turquesa, line width=0.75pt, -{Stealth[length=3.6pt]}]
    ($(#1)+({-90+#2}:#4)$) arc[start angle={-90+#2}, end angle={-90+#3}, radius=#4];}

% Pequeña pulsación: viñeta de la paginación y ornamento entre secciones
\newcommand{\ondita}[1][1]{%
  \tikz[baseline=-0.55ex, x=#1*1pt, y=#1*1pt]{%
    \draw[turquesa, line width=0.45pt, line join=round]
      plot[domain=0:18, samples=90, smooth, variable=\t]
      ({\t},{2.6*cos(2.1*\t r)*sin(\t*180/18)});
    \draw[azulprof!60, line width=0.3pt, densely dotted]
      plot[domain=0:18, samples=30, smooth, variable=\t] ({\t},{2.6*sin(\t*180/18)});}}


% =========================================================================
%        CABECERAS, PIE DE PÁGINA, SECCIONES E ÍNDICE
% =========================================================================
% Ornamento entre secciones: una pequeña pulsación entre dos filetes
\newcommand{\ornamento}{%
  \par\vspace{10pt}{\centering
  \tikz{\draw[azulclaro, line width=0.4pt] (-1.6,0) -- (-0.4,0);
        \draw[azulclaro, line width=0.4pt] (0.4,0) -- (1.6,0);
        \node at (0,0) {\ondita[1.25]};}\par}\vspace{6pt}}

\makeatletter
\def\ps@headings{%
  \let\@mkboth\markboth
  \def\@evenhead{\normalfont\footnotesize\color{grispizarra}\scshape
    \textls[60]{Laboratorio de Mecánica y Ondas}\hfil
    \textls[60]{Péndulos acoplados\,\textperiodcentered\,Modos normales}}%
  \let\@oddhead\@evenhead
  \def\@evenfoot{\hfil\normalfont\small\color{azulacero}%
    \ondita\enspace\thepage\enspace\ondita\hfil}%
  \let\@oddfoot\@evenfoot}
\def\ps@firstpage{\ps@plain\def\@oddhead{}\let\@evenhead\@oddhead
  \def\@oddfoot{\hfil\normalfont\small\color{azulacero}%
    \ondita\enspace\thepage\enspace\ondita\hfil}%
  \let\@evenfoot\@oddfoot}
\pagestyle{headings}

% Secciones: número en azul petróleo, título en versalitas espaciadas; debajo,
% un subtítulo en cursiva y un filete rematado en un pequeño círculo turquesa.
\def\section{\@startsection{section}{1}%
  \z@{1.3\linespacing\@plus.8\linespacing}{.25\linespacing}%
  {\normalfont\Large\color{azulprof}\scshape\lsstyle\raggedright}}
\let\@seccntformatorig\@seccntformat
\def\@seccntformat#1{%
  \ifnum\pdfstrcmp{#1}{section}=\z@
    {\color{azulacero}\upshape\csname the#1\endcsname}\hspace{0.55em}%
  \else\@seccntformatorig{#1}\fi}
\newcommand{\filete}{%
  \tikz[baseline]{\draw[azulclaro, line width=0.45pt] (0,0) -- (\linewidth-5pt,0);
    \draw[turquesa, line width=0.55pt, fill=white] (\linewidth-2.6pt,0) circle (2.3pt);}}
\DeclareRobustCommand{\tocsub}[1]{%
  \ \textperiodcentered\ {\itshape\color{grispizarra}#1}}
\newcommand{\subtitulo}[1]{%
  \par\nopagebreak\vspace{-2pt}%
  \noindent{\itshape\color{grispizarra}#1}\par\nopagebreak\vspace{1pt}%
  \noindent\filete\par\nopagebreak\vspace{0.55\baselineskip}}
% \seccion{Título}{Subtítulo}: sección numerada con subtítulo (índice incluido).
\let\@tocextra\@empty
\def\@tocwriteb#1#2#3{%
  \begingroup
    \def\@tocline##1##2##3##4##5##6{%
      \ifnum##1>\c@tocdepth
      \else \sbox\z@{##5\let\indentlabel\@tochangmeasure##6}\fi}%
    \csname l@#2\endcsname{#1{\csname#2name\endcsname}{\@secnumber}{}}%
  \endgroup
  \addcontentsline{toc}{#2}%
    {\protect#1{\csname#2name\endcsname}{\@secnumber}{#3\@tocextra}}%
  \global\let\@tocextra\@empty}
\newcommand{\seccion}[2]{%
  \gdef\@tocextra{\tocsub{#2}}%
  \section{#1}\subtitulo{#2}}

% Índice: número en azul petróleo, título en azul marino y subtítulo en cursiva
\def\l@section{\@tocline{1}{7pt}{0pt}{}{}}
\renewcommand{\tocsection}[3]{%
  \indentlabel{\@ifnotempty{#2}{{\color{azulacero}#2}\quad}}{\color{azulprof}#3}}
\def\l@subsection{\@tocline{2}{1pt}{1.55pc}{2.4pc}{\small}}
\renewcommand{\contentsnamefont}{\scshape\color{azulprof}\lsstyle}
\makeatother
\numberwithin{equation}{section}
\setcounter{tocdepth}{2}
\makeatletter
\def\@@and{y}
\renewcommand{\abstractname}{\color{azulacero}\lsstyle Resumen}
\makeatother

% =========================================================================
%               NOTACIÓN Y MACROS DE USO FRECUENTE
% =========================================================================
\newcommand{\dif}{\mathrm{d}}
\newcommand{\Tb}{T_{\mathrm{b}}}              % periodo de pulsación (envolvente)
\newcommand{\Tp}{T_{+}}                       % periodo de la portadora
\newcommand{\wz}{\omega_0}                    % frecuencia del modo en fase
\newcommand{\wu}{\omega_1}                    % frecuencia del modo en contrafase
\newcommand{\wm}{\omega_{-}}                  % frecuencia de modulación
\newcommand{\wpl}{\omega_{+}}                 % frecuencia de la portadora
\newcommand{\Tcoef}{t_{95,\nu}}               % factor de cobertura al 95 %
\newcommand{\ext}[1]{U(#1)}                   % incertidumbre ampliada
\newcommand{\mean}[1]{\overline{#1}}
% Dato pendiente de incorporar (se resalta con discreción)
\newcommand{\pend}[1]{\textcolor{grispizarra}{\colorbox{grisclaro}{\strut\,#1\,}}}
\newcommand{\pendiente}[1]{\textcolor{grispizarra}{\textit{[#1]}}}
% Hueco reservado para una gráfica o figura que aún no existe
\newcommand{\huecofigura}[2][5.2cm]{%
  \tikz{\draw[grispizarra!45, dashed, line width=0.5pt, rounded corners=2pt, fill=azulfondo]
          (0,0) rectangle (0.92\linewidth,#1);
        \node[text=grispizarra, font=\small\itshape, align=center, text width=0.8\linewidth]
          at (0.46\linewidth,#1/2) {#2};}}

% =========================================================================
%        ENTORNOS FORMALES (enunciados numerados, estilo artículo)
% =========================================================================
\usepackage{amsthm}

\newtheoremstyle{cienplano}{10pt}{10pt}{\itshape}{}%
  {\bfseries\color{azulprof}}{.}{0.6em}%
  {\thmname{#1}\thmnumber{ #2}\thmnote{\normalfont\upshape\ (#3)}}
\newtheoremstyle{ciendef}{10pt}{10pt}{\normalfont}{}%
  {\bfseries\color{azulprof}}{.}{0.6em}%
  {\thmname{#1}\thmnumber{ #2}\thmnote{\normalfont\upshape\ (#3)}}
\newtheoremstyle{cienobs}{8pt}{8pt}{\normalfont}{}%
  {\itshape\color{azulprof}}{.}{0.6em}%
  {\thmname{#1}\thmnumber{ #2}\thmnote{\normalfont\upshape\ (#3)}}

\theoremstyle{cienplano}
\newtheorem{proposicion}{Proposición}[section]
\newtheorem{teorema}[proposicion]{Teorema}
\newtheorem{resultado}[proposicion]{Resultado}
\newtheorem{conclusion}[proposicion]{Conclusión}
\theoremstyle{ciendef}
\newtheorem{definicion}[proposicion]{Definición}
\newtheorem{hipotesis}[proposicion]{Hipótesis}
\newtheorem{calculo}[proposicion]{Cálculo}
\newtheorem{estimacion}[proposicion]{Estimación lineal}
\theoremstyle{cienobs}
\newtheorem{observacion}[proposicion]{Observación}
\newtheorem*{demostracion}{Demostración}

% --- Referencia formalizada --------------------------------------------------
% Caja fina con fondo tenue y título sobre el borde. Cuerpo en cursiva y, al
% pie, la fuente completa con autor, título y número de cita.
% Uso: \begin{referencia}[título]{clave-bib} ... \end{referencia}
\newcounter{referencia}
\renewcommand{\thereferencia}{\theproposicion}
\newtcolorbox{referenciabox}[2]{%
  enhanced, sharp corners,
  colback=azulfondo, colframe=azulclaro, boxrule=0.4pt,
  left=10pt, right=10pt, top=7pt, bottom=6pt,
  before skip=14pt, after skip=12pt,
  attach boxed title to top left={yshift*=-\tcboxedtitleheight/2, xshift=8pt},
  boxed title style={colback=white, colframe=white, sharp corners, boxrule=0pt,
                     left=3pt, right=3pt, top=1pt, bottom=1pt},
  fonttitle=\bfseries\color{azulprof}, coltitle=azulprof,
  fontupper=\itshape,
  title={Referencia~\thereferencia\ifx\relax#1\relax\else\ \normalfont(#1)\fi},
  after upper={\par\vspace{5pt}{\color{azulclaro}\hrule height 0.4pt}\vspace{4pt}%
    {\upshape\footnotesize\color{grispizarra}\raggedleft
     \citeauthor{#2}, \citetitle{#2}~\cite{#2}.\par}}}
\newenvironment{referencia}[2][]
  {\par\stepcounter{proposicion}\refstepcounter{referencia}\begin{referenciabox}{#1}{#2}}
  {\end{referenciabox}}
\crefname{referencia}{referencia}{referencias}
\Crefname{referencia}{Referencia}{Referencias}

% --- Método (caja azulada para los fundamentos del tratamiento de datos) -----
\newcounter{metodo}
\renewcommand{\themetodo}{\theproposicion}
\newtcolorbox{metodobox}[1]{%
  enhanced, breakable, sharp corners,
  colback=azulfondo, colframe=azulclaro, boxrule=0.4pt,
  left=10pt, right=10pt, top=7pt, bottom=5pt,
  before skip=14pt, after skip=12pt,
  attach boxed title to top left={yshift*=-\tcboxedtitleheight/2, xshift=8pt},
  boxed title style={colback=white, colframe=white, sharp corners, boxrule=0pt,
                     left=3pt, right=3pt, top=1pt, bottom=1pt},
  fonttitle=\bfseries\color{azulprof}, coltitle=azulprof,
  title={Método~\themetodo\ifx\relax#1\relax\else\ \normalfont(#1)\fi}}
\newenvironment{metodo}[1][]
  {\par\stepcounter{proposicion}\refstepcounter{metodo}\begin{metodobox}{#1}}
  {\end{metodobox}}
\crefname{metodo}{método}{métodos}
\Crefname{metodo}{Método}{Métodos}

% --- Resultado destacado ----------------------------------------------------
\newtcolorbox{destacado}[1][]{%
  enhanced, sharp corners, frame hidden, colback=azulfondo,
  borderline west={1.6pt}{0pt}{turquesa},
  left=10pt, right=8pt, top=2pt, bottom=4pt,
  before skip=8pt, after skip=10pt,
  before upper={\ifx\relax#1\relax\else{\scshape\small\color{azulacero}#1}\par\vspace{-4pt}\fi}}

% --- Tarea pendiente para el compañero ---------------------------------------
% Caja discontinua para las partes que aún hay que redactar o calcular.
\newtcolorbox{pendientebox}[1][]{%
  enhanced, breakable, sharp corners, colback=white, colframe=grispizarra!50,
  boxrule=0.5pt, borderline={0.5pt}{0pt}{grispizarra!50, dashed}, frame hidden,
  left=10pt, right=10pt, top=6pt, bottom=6pt, before skip=10pt, after skip=10pt,
  fontupper=\small\color{grispizarra},
  before upper={{\scshape\color{azulacero}Pendiente\ifx\relax#1\relax\else\ \textperiodcentered\ #1\fi}\par\vspace{2pt}}}

\crefname{proposicion}{proposición}{proposiciones}
\Crefname{proposicion}{Proposición}{Proposiciones}
\crefname{teorema}{teorema}{teoremas}
\Crefname{teorema}{Teorema}{Teoremas}
\crefname{definicion}{definición}{definiciones}
\Crefname{definicion}{Definición}{Definiciones}
\crefname{resultado}{resultado}{resultados}
\Crefname{resultado}{Resultado}{Resultados}
\crefname{conclusion}{conclusión}{conclusiones}
\crefname{hipotesis}{hipótesis}{hipótesis}
\crefname{calculo}{cálculo}{cálculos}
\crefname{estimacion}{estimación}{estimaciones}
\crefname{observacion}{observación}{observaciones}
\Crefname{observacion}{Observación}{Observaciones}
\crefname{section}{sección}{secciones}
\Crefname{section}{Sección}{Secciones}
\crefname{subsection}{sección}{secciones}
\Crefname{subsection}{Sección}{Secciones}
\crefname{subsubsection}{sección}{secciones}
\crefname{table}{tabla}{tablas}
\Crefname{table}{Tabla}{Tablas}
\crefname{figure}{figura}{figuras}
\Crefname{figure}{Figura}{Figuras}
\crefname{equation}{ecuación}{ecuaciones}
\Crefname{equation}{Ecuación}{Ecuaciones}
\crefformat{equation}{ecuación~(#2#1#3)}
\Crefformat{equation}{Ecuación~(#2#1#3)}
\crefrangeformat{equation}{ecuaciones~(#3#1#4)--(#5#2#6)}
\crefmultiformat{equation}{ecuaciones~(#2#1#3)}{ y~(#2#1#3)}{, (#2#1#3)}{ y~(#2#1#3)}

% --- Estilo de gráficas de datos ----------------------------------------------
\pgfplotsset{
  datos/.style={
    tick label style={font=\scriptsize}, label style={font=\small},
    axis lines=left, axis line style={grispizarra},
    grid=major, grid style={grisclaro, line width=0.4pt},
    every tick/.style={grispizarra},
    legend style={font=\scriptsize, draw=none, fill=white, fill opacity=0.9,
                  text opacity=1, cells={anchor=west}},
    legend cell align=left,
    /pgf/number format/use comma, /pgf/number format/1000 sep={\,},
    scaled ticks=false,
  },
  senal/.style={azulprof, line width=0.55pt},
  senalB/.style={turquesa!90!azulprof, line width=0.55pt},
  envolvente/.style={azulacero!70, line width=0.6pt, densely dashed},
  puntoA/.style={only marks, mark=*, mark size=2.1pt, azulprof,
    mark options={fill=azulprof, draw=white, line width=0.3pt}},
  puntoB/.style={only marks, mark=square*, mark size=1.9pt, azulacero,
    mark options={fill=azulclaro, draw=azulacero, line width=0.5pt}},
  leyenda arriba/.style={legend columns=-1,
    legend style={at={(0.5,1.03)}, anchor=south,
      /tikz/every even column/.append style={column sep=9pt}}},
  teorica/.style={acento, line width=0.9pt, densely dashed},
}

% Referencia a un panel de figura: «figura 2\,(a)», enlazada en bloque
\newcommand{\figpanel}[2]{\hyperref[#1]{figura~\ref*{#1}\,\textup{(#2)}}}

% Pies de figura y tabla
\usepackage[font=small, labelfont={sc,color=azulprof}, labelsep=period,
            margin=0.6cm, justification=justified]{caption}
\usepackage{subcaption}
\captionsetup[sub]{font=small, labelfont={sc,color=azulprof}, labelformat=parens,
                   margin=0pt, justification=centering}

% Fotografías con un marco fino
\newcommand{\foto}[2][\linewidth]{%
  \tcbox[enhanced, sharp corners, colback=white, colframe=grisclaro, boxrule=0.5pt,
         boxsep=0pt, left=0pt, right=0pt, top=0pt, bottom=0pt, nobeforeafter, valign=center]{%
    \includegraphics[width=\dimexpr#1-1pt\relax]{#2}}}

% =========================================================================
%                     VARIABLES PARA LA PORTADA
% =========================================================================
\newcommand{\Universidad}{Universidad de Salamanca}
\newcommand{\Facultad}{Facultad de Ciencias}
\newcommand{\Titulacion}{Grado en Física}
\newcommand{\Asignatura}{Laboratorio de Mecánica y Ondas}
\newcommand{\NumPractica}{Práctica 3}
\newcommand{\TituloPractica}{Péndulos acoplados}
\newcommand{\SubtituloPractica}{Modos normales de oscilación y pulsaciones}
\newcommand{\tipotrabajo}{Informe de laboratorio}
\newcommand{\miNombre}{David Gómez García \\ Daniel Martín Guerrero}
\newcommand{\miTutor}{Mario Amado \\ Ivan de Martino \\ Manuel Pino}
\newcommand{\miUbicacion}{Salamanca}

% =========================================================================
%                     DATOS DEL ARTÍCULO (cabecera)
% =========================================================================
\title{Péndulos acoplados}
\author{David Gómez García}
\author{Daniel Martín Guerrero}
\address{Grado en Física, Facultad de Ciencias, Universidad de Salamanca,
         Salamanca, España}
\date{}

% Valores calculados a partir de la hoja de cálculo (scripts/tablas.py)
% Generado por scripts/tablas.py a partir de Practica3_pendulos_acoplados.xlsx. No editar a mano.

\newcommand{\valTcero}{\num{1.9960}}
\newcommand{\valWcero}{\num{3.1479}}
\newcommand{\valWceroDos}{\num{9.9092}}
\newcommand{\valI}{\num{0.9900}}


% =========================================================================
%                                DOCUMENTO
% =========================================================================
\begin{document}
% --- PORTADA ----------------------------------------------------------------
\pagenumbering{alph}
%%%%%%%%%%%%%%%%%%%%%%%%%%%%%%%%%%%%%%%%%%%%%%%%%%%%%%%%%%%%%%%%%%%%%%%%
% Portada — estilo artículo científico, motivo: péndulos acoplados
% Misma estructura que la del informe de Steiner (cabecera con logotipos,
% motivo gráfico, franja inferior, bloque de título y pie con autores),
% ahora en la gama de azul marino, petróleo y turquesa. El motivo es el
% montaje del laboratorio dibujado de frente, en un instante del modo
% pulsante; debajo, la señal de los dos péndulos con su envolvente.
% Requiere (definidos en main.tex): colores azulprof, azulacero, turquesa,
% azulclaro, grispizarra, grisclaro, azulfondo; macros \Universidad,
% \Facultad, \Titulacion, \Asignatura, \NumPractica, \TituloPractica,
% \SubtituloPractica, \tipotrabajo, \miNombre, \miTutor, \miUbicacion;
% y figuras/pendulos.tex.
%%%%%%%%%%%%%%%%%%%%%%%%%%%%%%%%%%%%%%%%%%%%%%%%%%%%%%%%%%%%%%%%%%%%%%%%

\begin{titlepage}
\thispagestyle{empty}
\pagecolor{white}

\begin{tikzpicture}[remember picture, overlay]
  \coordinate (O) at (current page.south west);

  % ----------------------------------------------------------------
  % Cabecera: logotipos y filete
  % ----------------------------------------------------------------
  \node[anchor=north west, inner sep=0] at ($(O)+(2.2cm,27.9cm)$)
    {\includegraphics[height=1.42cm]{figuras/logo-usal.png}};
  \node[anchor=north east, inner sep=0] at ($(O)+(18.8cm,27.9cm)$)
    {\includegraphics[height=1.42cm]{figuras/logo-ciencias.png}};
  \draw[grispizarra!60, line width=0.4pt]
    ($(O)+(2.2cm,26.0cm)$) -- ($(O)+(18.8cm,26.0cm)$);
  \node[anchor=north west, inner sep=0, text=grispizarra, font=\scshape\small]
    at ($(O)+(2.2cm,25.75cm)$)
    {\textls[80]{\Universidad\ \textbar\ \Facultad}};
  \node[anchor=north east, inner sep=0, text=grispizarra, font=\scshape\small]
    at ($(O)+(18.8cm,25.75cm)$)
    {\textls[80]{\Titulacion}};

  % ----------------------------------------------------------------
  % Motivo gráfico: los dos péndulos acoplados, vistos de frente, en un
  % instante del modo pulsante (el izquierdo casi en reposo, el derecho
  % con toda la amplitud). Arcos punteados: trayectoria de las lentejas.
  % ----------------------------------------------------------------
  \begin{scope}[shift={($(O)+(10.5cm,23.75cm)$)}, x=1cm, y=1cm]
    \def\Lp{6.55}\def\lp{2.95}\def\dsep{5.6}
    % barra de soporte, muy discreta
    \draw[grisclaro, line width=2.2pt, line cap=round] (-4.6,0.95) -- (4.6,0.95);
    % trayectorias de las lentejas y escala angular tenue
    \foreach \cx in {-2.8,2.8}{
      \draw[azulclaro, line width=0.5pt, densely dotted]
        ($(\cx,0)+(-104:\Lp)$) arc[start angle=-104, end angle=-76, radius=\Lp];
      \foreach \a in {-100,-95,...,-80}{
        \draw[azulclaro!80, line width=0.3pt]
          ($(\cx,0)+(\a:\Lp+0.42)$) -- ($(\cx,0)+(\a:\Lp+0.62)$);}}
    % el par de péndulos
    \parpendulos{-3}{12}{\lp}{\Lp}{\dsep}{1.2}
    % ángulos
    \angulo{P2}{12}{2.0}{$\varphi_2$}
    \angulo{P1}{-3}{2.0}{$\varphi_1$}
    % longitud de acoplamiento
    \draw[grispizarra, line width=0.45pt, |<->|]
      ($(P2)+(0.55,0)$) -- ($(S2)+(0.55,0)$)
      node[midway, right=1pt, text=azulprof, fill=white, inner sep=1.5pt] {$\ell$};
    % longitud del péndulo
    \draw[grispizarra!70, line width=0.4pt, |<->|]
      ($(P1)+(-0.75,0)$) -- ($(P1)+(-0.75,-\Lp)$)
      node[pos=0.62, left=1pt, text=azulprof, fill=white, inner sep=1.5pt] {$L$};
    % rótulos discretos
    \draw[grispizarra!70, line width=0.35pt] ($(S1)!0.5!(S2)+(0,-0.12)$) -- ++(0.9,-1.25);
    \fill[azulacero] ($(S1)!0.5!(S2)+(0,-0.12)$) circle (1.3pt);
    \node[anchor=west, text=grispizarra, font=\itshape\small]
      at ($(S1)!0.5!(S2)+(0.95,-1.32)$) {muelle, $k$};
    \draw[grispizarra!70, line width=0.35pt] ($(B2)+(0.48,0.1)$) -- ++(1.1,0.55);
    \node[anchor=west, text=grispizarra, font=\itshape\small]
      at ($(B2)+(1.62,0.68)$) {lenteja, $m$};
    \draw[grispizarra!70, line width=0.35pt] ($(P1)+(-0.85,0.45)$) -- ++(-1.0,0.45);
    \node[anchor=east, text=grispizarra, font=\itshape\small]
      at ($(P1)+(-1.9,0.92)$) {eje de giro y sensor};
  \end{scope}

  % ----------------------------------------------------------------
  % Franja: la pulsación. Los dos péndulos se pasan la energía; cada
  % trazo es una portadora rápida dentro de una envolvente lenta, y las
  % envolventes están desfasadas un cuarto de ciclo.
  % ----------------------------------------------------------------
  \begin{scope}[shift={($(O)+(0,13.55cm)$)}, x=1cm, y=1cm]
    \def\xa{2.2}\def\xb{18.8}\def\wm{0.2838}\def\wp{6.6}
    \draw[grispizarra!35, line width=0.3pt] (\xa,0.62) -- (\xb,0.62);
    \draw[grispizarra!35, line width=0.3pt] (\xa,-0.62) -- (\xb,-0.62);
    % péndulo 2 (arriba): cos(w- t) cos(w+ t)
    \draw[azulacero!55, line width=0.35pt, densely dashed, domain=\xa:\xb, samples=120]
      plot (\x,{0.62+0.48*abs(cos(\wm*(\x-\xa) r))});
    \draw[azulprof, line width=0.5pt, domain=\xa:\xb, samples=1100]
      plot (\x,{0.62+0.48*cos(\wm*(\x-\xa) r)*cos(\wp*(\x-\xa) r)});
    % péndulo 1 (abajo): sin(w- t) sin(w+ t)
    \draw[turquesa!60, line width=0.35pt, densely dashed, domain=\xa:\xb, samples=120]
      plot (\x,{-0.62+0.48*abs(sin(\wm*(\x-\xa) r))});
    \draw[turquesa!90!azulprof, line width=0.5pt, domain=\xa:\xb, samples=1100]
      plot (\x,{-0.62+0.48*sin(\wm*(\x-\xa) r)*sin(\wp*(\x-\xa) r)});
    \node[anchor=east, text=azulprof, font=\footnotesize] at (\xa-0.08,0.62) {$\varphi_2$};
    \node[anchor=east, text=turquesa!90!azulprof, font=\footnotesize] at (\xa-0.08,-0.62) {$\varphi_1$};
    % periodo de pulsación sobre la envolvente superior
    \draw[grispizarra, line width=0.4pt, |<->|]
      (\xa,1.32) -- ({\xa+pi/\wm},1.32)
      node[midway, above=0pt, text=azulprof, font=\footnotesize, fill=white, inner sep=1.2pt] {$T_{\mathrm{b}}=\pi/\omega_{-}$};
    \node[anchor=north, text=grispizarra, font=\scshape\scriptsize] at (10.5,-1.32)
      {\textls[160]{la energía pasa de un péndulo al otro}};
  \end{scope}

  % ----------------------------------------------------------------
  % Bloque del título
  % ----------------------------------------------------------------
  \node[anchor=north west, inner sep=0, text width=16.6cm, align=left]
    at ($(O)+(2.2cm,11.35cm)$)
    {{\color{azulacero}\scshape\large\textls[120]{\Asignatura\ \textbar\ \NumPractica}}\\[12pt]
     {\color{azulprof}\fontsize{34}{39}\selectfont\bfseries\spaceskip=0.28em\relax\TituloPractica}\\[10pt]
     {\color{azulprof}\itshape\Large\SubtituloPractica}\\[12pt]
     {\color{turquesa}\rule{2.5cm}{1.4pt}}\\[9pt]
     {\color{grispizarra}\large
      $\omega_1^{2}=\omega_0^{2}+\dfrac{2k}{I}\,\ell^{2}$
      \quad\textperiodcentered\quad
      $\omega_{\pm}=\dfrac{\omega_1\pm\omega_0}{2}$}};

  % ----------------------------------------------------------------
  % Pie: tipo de trabajo, autores, tutores, lugar y fecha
  % ----------------------------------------------------------------
  \draw[grispizarra!60, line width=0.4pt]
    ($(O)+(2.2cm,4.6cm)$) -- ($(O)+(18.8cm,4.6cm)$);

  \node[anchor=north west, inner sep=0, text width=4.6cm, align=left]
    at ($(O)+(2.2cm,4.25cm)$)
    {{\scshape\small\color{azulacero}\textls[80]{Documento}}\\[3pt]
     {\large\bfseries\color{azulprof}\tipotrabajo}};

  \node[anchor=north west, inner sep=0, text width=5.6cm, align=left]
    at ($(O)+(7.2cm,4.25cm)$)
    {{\scshape\small\color{azulacero}\textls[80]{Autores}}\\[3pt]
     {\large\color{black!85}\miNombre}};

  \node[anchor=north west, inner sep=0, text width=5cm, align=left]
    at ($(O)+(13.8cm,4.25cm)$)
    {{\scshape\small\color{azulacero}\textls[80]{Tutores}}\\[3pt]
     {\large\color{black!85}\miTutor}};

  \node[anchor=south east, inner sep=0, text=grispizarra, font=\scshape\small]
    at ($(O)+(18.8cm,1.2cm)$)
    {\textls[100]{\miUbicacion, \number\year}};
  \fill[azulprof] ($(O)+(0,0)$) rectangle ($(O)+(21cm,0.45cm)$);
  \fill[azulacero] ($(O)+(0,0)$) rectangle ($(O)+(9cm,0.45cm)$);
  \fill[turquesa] ($(O)+(0,0)$) rectangle ($(O)+(5cm,0.45cm)$);
  \fill[azulclaro] ($(O)+(0,0)$) rectangle ($(O)+(2cm,0.45cm)$);
\end{tikzpicture}

\end{titlepage}

\clearpage
\pagenumbering{arabic}
\setcounter{page}{1}

\begin{abstract}
Estudiamos dos péndulos físicos iguales unidos por un muelle sujeto a una
distancia $\ell$ de sus ejes de giro. En el régimen de pequeñas oscilaciones
el sistema tiene dos modos normales: uno en fase, de frecuencia $\wz$, en el
que el muelle no llega a deformarse, y otro en contrafase, de frecuencia
$\wu$, que crece con $\ell$. Su superposición da lugar al modo pulsante, en
el que la energía pasa de un péndulo al otro de forma periódica. Registramos
los tres modos para cinco longitudes de acoplamiento, entre
\SIlist{25;65}{\centi\meter}, con sensores de ángulo, la unidad Cobra-3 y el
programa PHYWE Measure, cuyas herramientas de suavizado y de análisis de
picos nos dan los instantes de los máximos y, de ellos, los periodos. Del
modo en fase obtenemos $\wz=\valWcero\ \si{\radian\per\second}$ y, con él,
el momento de inercia de cada péndulo, $I=\valI\ \si{\kilogram\meter\squared}$.
Después ajustamos linealmente $\wu^{2}$, $\wm$ y $\wpl$ frente a $\ell^{2}$ y,
de cada pendiente, deducimos la constante elástica del muelle,
$k=$\,\pend{$k_{1}$}, \pend{$k_{2}$} y \pend{$k_{3}$}, cuyos valores
comparamos para discutir qué método es más preciso. Los resultados se
expresan con su incertidumbre típica y ampliada al \SI{95}{\percent},
siguiendo \emph{Notes on Uncertainty in Measurements}.
\end{abstract}

\maketitle
\tableofcontents

% =========================================================================
\newpage
\seccion{Introducción}{Dos osciladores que se comunican}\label{sec:introduccion}
% =========================================================================
Un péndulo aislado es probablemente el oscilador más familiar: apartado un
poco de la vertical, vuelve a ella y la rebasa una y otra vez con un periodo
fijo. Basta unir dos péndulos con un muelle para que la historia cambie. El
muelle transmite fuerza, y con ella energía, de uno a otro, y el movimiento
de cada péndulo deja de ser una oscilación sencilla: crece, se apaga y vuelve
a crecer mientras su compañero hace lo contrario.

Ese movimiento, que a primera vista parece desordenado, tiene una estructura
muy precisa. Existen unas configuraciones especiales, los \emph{modos
normales}, en las que los dos péndulos oscilan con la misma frecuencia y una
relación de fase fija, y cualquier otro movimiento es una superposición de
ellos~\cite{FIV,FRENCH}. La idea va mucho más allá de dos péndulos: es la
misma que describe las vibraciones de una molécula, las de una red cristalina
o las de dos circuitos eléctricos acoplados, y es la puerta de entrada al
estudio de las ondas~\cite{FIV,CRAWFORD}.

En esta práctica la intensidad del acoplamiento se regula sin tocar los
péndulos: basta desplazar el muelle a lo largo de las varillas. Así podemos
separar lo que es propio de cada péndulo, su frecuencia natural $\wz$, de lo
que introduce el muelle, su constante elástica $k$~\cite{GUION-P3}.

\subsection{Objetivos}\label{sec:objetivos}
\begin{itemize}
  \item Generar y registrar los tres modos fundamentales de oscilación (en
        fase, en contrafase y pulsante) para cinco longitudes de
        acoplamiento $\ell$.
  \item Determinar, a partir del modo en fase, la frecuencia natural $\wz$
        de los péndulos y su momento de inercia $I$.
  \item Representar $\wu^{2}$, $\wm$ y $\wpl$ frente a $\ell^{2}$, ajustar
        rectas por mínimos cuadrados y deducir de cada pendiente la
        constante $k$ del muelle.
  \item Comparar los tres valores de $k$ y discutir qué método proporciona el
        resultado más preciso y por qué.
\end{itemize}
Para ello nos apoyaremos en las \cref{eq:inercia,eq:lineales,eq:k} de la
\cref{sec:fundamentos}.

\subsection{Material empleado}\label{sec:material}
Hemos utilizado el siguiente instrumental (\cref{fig:fotos-montaje}):
\begin{itemize}
  \item Dos péndulos de gravitación iguales, de masa $m=\SI{1}{\kilogram}$ y
        longitud $L=\SI{1}{\meter}$~\cite{GUION-P3}, formados por una varilla
        con una masa lenticular en su extremo; sus ejes de giro llevan un
        sensor de ángulo.
  \item Muelle helicoidal de acoplamiento, sujeto a cada varilla con una
        pinza deslizante.
  \item Fuente de alimentación y cables de conexión.
  \item Unidad de medida Cobra-3 y ordenador con el programa PHYWE
        Measure~\cite{PHYWE}.
  \item Cinta métrica, para la longitud de acoplamiento.
\end{itemize}

\subsection{Montaje experimental}\label{sec:montaje}
Cada péndulo cuelga de un cabezal triangular fijado al borde de la mesa con
una mordaza (\cref{fig:fotos-montaje}). En el eje de giro, un sensor de
ángulo convierte la desviación de la varilla en una tensión que la unidad
Cobra-3 registra en dos canales analógicos, uno por péndulo, de modo que el
programa dibuja en tiempo real dos curvas tensión--tiempo proporcionales a
$\varphi_1(t)$ y $\varphi_2(t)$. El muelle une las dos varillas a la misma
altura; la distancia $\ell$ entre los ejes de giro y los puntos de sujeción
es la \emph{longitud de acoplamiento}, que medimos con la cinta métrica y
que cambiamos deslizando a la vez las dos pinzas (\cref{fig:esquema}).

\begin{figure}[ht]
\centering
\begin{minipage}[b]{0.355\textwidth}\centering
  \foto{figuras/montaje-foto.jpg}\\[3pt]
  {\small\scshape (a)\; montaje completo}
\end{minipage}\hfill
\begin{minipage}[b]{0.6\textwidth}\centering
  \foto{figuras/montaje-muelle.jpg}\\[3pt]
  {\small\scshape (b)\; muelle de acoplamiento}\\[10pt]
  \begin{minipage}{\linewidth}\small
  La unidad Cobra-3 (al fondo, a la izquierda) recoge la señal de los dos
  sensores; la fuente de alimentación (a la derecha) los alimenta. Las
  masas lenticulares, de color naranja, cuelgan a $L\approx\SI{1}{\meter}$
  de los ejes, y el muelle puede situarse a cualquier altura de las varillas.
  \end{minipage}
\end{minipage}
\caption{El montaje en el laboratorio. (a)~Los dos péndulos, con los
cabezales triangulares que alojan los ejes de giro y los sensores. (b)~Detalle
del muelle helicoidal, sujeto a cada varilla mediante una pinza a la
distancia $\ell$ del eje.}
\label{fig:fotos-montaje}
\end{figure}

\begin{figure}[ht]
\centering
\begin{tikzpicture}[>=Stealth, x=0.78cm, y=0.78cm, font=\small]
  \def\Lm{7.2}\def\lm{3.2}\def\dm{5.4}
  \parpendulos{7}{-5}{\lm}{\Lm}{\dm}{1}
  \angulo{P1}{7}{2.0}{$\varphi_1$}
  \angulo{P2}{-5}{2.0}{$\varphi_2$}
  % cotas
  \draw[grispizarra, line width=0.45pt, |<->|]
    ($(P2)+(0.95,0)$) -- ($(P2)+(0.95,-\lm)$)
    node[midway, right=1pt, text=azulprof] {$\ell$};
  \draw[grispizarra, line width=0.45pt, |<->|]
    ($(P1)+(-1.0,0)$) -- ($(P1)+(-1.0,-\Lm)$)
    node[pos=0.6, left=1pt, text=azulprof] {$L$};
  \draw[grispizarra!50, densely dotted, line width=0.35pt] (S2) -- ($(P2)+(1.1,-\lm)$);
  % rótulos
  \node[text=azulprof, anchor=west] at ($(B2)+(0.55,0)$) {$m$};
  \node[text=grispizarra, font=\itshape\footnotesize, anchor=south] at ($(S1)!0.5!(S2)+(0,0.22)$) {muelle, $k$};
  \node[text=grispizarra, font=\itshape\footnotesize, anchor=south] at ($(P1)!0.5!(P2)+(0,0.25)$) {ejes de giro (sensores)};
  % fuerzas del muelle (péndulos acercándose: el muelle está comprimido)
  \draw[acento, line width=0.8pt, ->] (S1) -- ++(-0.95,0) node[above, font=\footnotesize] {$F$};
  \draw[acento, line width=0.8pt, ->] (S2) -- ++(0.95,0) node[above, font=\footnotesize] {$F$};
  % peso
  \draw[azulacero, line width=0.8pt, ->] ($(B1)+(0,-0.42)$) -- ++(0,-1.2) node[right, font=\footnotesize] {$m\vec{g}$};
\end{tikzpicture}
\caption{Esquema del montaje, de frente y sin escala. Cada péndulo, de masa
$m$ y longitud $L$, gira alrededor de un eje horizontal; $\varphi_1$ y
$\varphi_2$ son sus desviaciones respecto de la vertical (discontinua),
positivas hacia la derecha. El muelle, de constante $k$, une las varillas a
una distancia $\ell$ de los ejes. En el instante dibujado los péndulos se
acercan, el muelle se comprime y empuja a cada uno hacia fuera con una fuerza
$F=k\ell(\varphi_1-\varphi_2)$.}
\label{fig:esquema}
\end{figure}

\ornamento

% =========================================================================
\newpage
\seccion{Fundamentos teóricos}{De dos ecuaciones acopladas a dos osciladores independientes}\label{sec:fundamentos}
% =========================================================================
\subsection{Ecuaciones del movimiento}\label{sec:ecuaciones}
Cada péndulo es un sólido rígido que gira alrededor de un eje horizontal fijo,
con momento de inercia $I$ respecto a él. Sobre el péndulo $j$ actúan dos
momentos. El primero es el del peso, aplicado a la distancia $L$ del eje,
\begin{equation}\label{eq:peso}
  \tau_{g}^{(j)}=-mgL\sin\varphi_j\simeq-mgL\,\varphi_j ,
  \qquad j=1,2,
\end{equation}
donde la aproximación vale para ángulos pequeños, el régimen en el que
trabajamos. El segundo es el del muelle. Si las varillas giran $\varphi_1$ y
$\varphi_2$, los puntos de sujeción se desplazan horizontalmente
$\ell\varphi_1$ y $\ell\varphi_2$, de modo que el muelle cambia su longitud en
$\ell(\varphi_2-\varphi_1)$ y ejerce una fuerza $k\ell(\varphi_2-\varphi_1)$ con
brazo $\ell$. Sus momentos sobre los dos péndulos son iguales y opuestos:
\begin{equation}\label{eq:muelle}
  \tau_{k}^{(1)}=k\ell^{2}(\varphi_2-\varphi_1),
  \qquad
  \tau_{k}^{(2)}=-k\ell^{2}(\varphi_2-\varphi_1).
\end{equation}
Llevando ambos a la ecuación fundamental de la dinámica de rotación,
$\tau=I\ddot\varphi$~\cite{GUION-P3}, obtenemos
\begin{equation}\label{eq:acopladas}
\begin{aligned}
  I\ddot\varphi_1&=-mgL\,\varphi_1+k\ell^{2}(\varphi_2-\varphi_1),\\
  I\ddot\varphi_2&=-mgL\,\varphi_2-k\ell^{2}(\varphi_2-\varphi_1).
\end{aligned}
\end{equation}

\begin{definicion}[Frecuencias características]\label{def:frecuencias}
Llamamos \emph{frecuencia natural} del péndulo aislado y \emph{frecuencia de
acoplamiento} a
\begin{equation}\label{eq:defs}
  \wz^{2}=\frac{mgL}{I},
  \qquad
  \Omega^{2}=\frac{k\ell^{2}}{I}.
\end{equation}
La primera solo depende del péndulo; la segunda, del muelle y de dónde lo
sujetamos.
\end{definicion}

Con ellas, el sistema de la \cref{eq:acopladas} toma una forma más limpia, que conviene
escribir también matricialmente, como en los apuntes de Física~IV~\cite{FIV}:
\begin{equation}\label{eq:matricial}
  \begin{pmatrix}\ddot\varphi_1\\ \ddot\varphi_2\end{pmatrix}
  =-\begin{pmatrix}\wz^{2}+\Omega^{2} & -\Omega^{2}\\
                   -\Omega^{2} & \wz^{2}+\Omega^{2}\end{pmatrix}
   \begin{pmatrix}\varphi_1\\ \varphi_2\end{pmatrix}.
\end{equation}
Los términos fuera de la diagonal son el acoplamiento: la aceleración de cada
péndulo depende también de la posición del otro.

\subsection{Modos normales}\label{sec:modos}
La matriz de la \cref{eq:matricial} es simétrica y sus dos elementos
diagonales son iguales, así que el sistema se desacopla con solo sumar y
restar las ecuaciones~\cite{FIV}. Definimos las \emph{coordenadas normales}
\begin{equation}\label{eq:coordnormales}
  q_0=\varphi_1+\varphi_2,
  \qquad
  q_1=\varphi_1-\varphi_2,
\end{equation}
que obedecen, cada una por su cuenta, la ecuación de un oscilador armónico:
\begin{equation}\label{eq:desacopladas}
  \ddot q_0+\wz^{2}\,q_0=0,
  \qquad
  \ddot q_1+\wu^{2}\,q_1=0 .
\end{equation}

\begin{resultado}[Frecuencias de los modos normales]\label{res:modos}
Los dos péndulos acoplados tienen dos modos normales: uno simétrico, de
frecuencia $\wz$, y otro antisimétrico, de frecuencia
\begin{equation}\label{eq:modos}
  \wu=\sqrt{\wz^{2}+2\Omega^{2}}
     =\wz\sqrt{1+\frac{2k\ell^{2}}{mgL}}\;>\;\wz .
\end{equation}
\end{resultado}

En el modo simétrico ($q_1=0$) los péndulos van a la par y el muelle nunca se
deforma: por eso su frecuencia es la del péndulo aislado. En el antisimétrico
($q_0=0$) el muelle trabaja con toda su fuerza y añade rigidez al sistema, lo
que eleva la frecuencia tanto más cuanto más abajo lo sujetamos. Deshaciendo
el cambio de la \cref{eq:coordnormales}, la solución general es
\begin{equation}\label{eq:general}
\begin{aligned}
  \varphi_1(t)&=\tfrac12\bigl[A_0\cos(\wz t+\alpha_0)+A_1\cos(\wu t+\alpha_1)\bigr],\\
  \varphi_2(t)&=\tfrac12\bigl[A_0\cos(\wz t+\alpha_0)-A_1\cos(\wu t+\alpha_1)\bigr],
\end{aligned}
\end{equation}
con las amplitudes $A_0$, $A_1$ y las fases $\alpha_0$, $\alpha_1$ fijadas
por las condiciones iniciales. Es exactamente lo que recogen los apuntes:

\begin{referencia}[Superposición de modos normales]{FIV}\label{ref:superposicion}
Cualquier oscilación arbitraria que se le imprima al sistema de osciladores
acoplados puede describirse como una combinación lineal de los modos
normales del sistema.
\end{referencia}

\subsection{Los tres modos fundamentales}\label{sec:tresmodos}
Partiendo siempre del reposo ($\dot\varphi_1(0)=\dot\varphi_2(0)=0$, de modo
que $\alpha_0=\alpha_1=0$), la \cref{eq:general} se particulariza en tres
casos sencillos, que son los que generamos en el laboratorio
(\cref{fig:regimenes}).

\begin{figure}[ht]
\centering
\begin{tikzpicture}[x=0.8cm, y=0.8cm, >=Stealth]
  \def\Lr{2.15}\def\lr{0.95}\def\dr{1.4}\def\sr{0.42}\def\Rm{2.5}
  % \panel{x}{phi1}{phi2}{fantasma1}{fantasma2}{rótulo de tiempo}
  % (si un péndulo está en reposo, su fantasma es "no")
  \newcommand{\panel}[6]{%
    \begin{scope}[shift={(#1,0)}]
      \def\fa{#4}\def\fb{#5}\def\no{no}%
      \ifx\fa\no\else
        \fantasma{-0.7,0}{#4}{\Lr}{\sr}
        \movimiento{-0.7,0}{0.55*#2}{0.55*#4}{\Rm}\fi
      \ifx\fb\no\else
        \fantasma{0.7,0}{#5}{\Lr}{\sr}
        \movimiento{0.7,0}{0.55*#3}{0.55*#5}{\Rm}\fi
      \parpendulos{#2}{#3}{\lr}{\Lr}{\dr}{\sr}
      \node[text=grispizarra, font=\footnotesize] at (0,-3.05) {#6};
    \end{scope}}
  % (a) en fase
  \panel{0}{13}{13}{-13}{-13}{$t$}
  \panel{3.0}{-13}{-13}{13}{13}{$t+T_0/2$}
  % (b) contrafase
  \panel{6.5}{-13}{13}{13}{-13}{$t$}
  \panel{9.5}{13}{-13}{-13}{13}{$t+T_1/2$}
  % (c) pulsante: la oscilación pasa del péndulo 2 al 1
  \panel{13.0}{0}{13}{no}{-13}{$t$}
  \panel{16.0}{-13}{0}{13}{no}{$t+\Tb/2$}
  % títulos y separadores
  \foreach \x/\t/\s in {1.5/{(a)\; en fase}/{$\varphi_1=\varphi_2$, frecuencia $\wz$},
                         8.0/{(b)\; contrafase}/{$\varphi_1=-\varphi_2$, frecuencia $\wu$},
                         14.5/{(c)\; pulsante}/{la energía pasa de un péndulo al otro}}{
    \node[anchor=south, align=center, inner sep=1pt] at (\x,0.62)
      {{\scshape\small\color{azulprof}\t}\\[-1pt]{\footnotesize\color{grispizarra}\s}};}
  \foreach \x in {4.75,11.25}{
    \draw[grisclaro, line width=0.6pt] (\x,1.35) -- (\x,-3.25);}
\end{tikzpicture}
\caption{Los tres modos fundamentales, cada uno en dos instantes separados
medio periodo. En trazo oscuro, la posición de los péndulos en ese instante;
en claro, el extremo opuesto de su oscilación, hacia el que se mueven
(flechas). (a)~En fase, el muelle conserva su longitud. (b)~En contrafase, se
estira y se comprime alternativamente. (c)~En el modo pulsante, la oscilación
pasa de un péndulo al otro en medio periodo de pulsación, $\Tb/2$, mientras
el compañero queda en reposo.}
\label{fig:regimenes}
\end{figure}

\subsubsection*{Oscilación en fase}
Si soltamos los dos péndulos con la misma desviación, $\varphi_1(0)=\varphi_2(0)$,
entonces $A_1=0$ y
\begin{equation}\label{eq:fase}
  \varphi_1(t)=\varphi_2(t)=\varphi_1(0)\cos(\wz t).
\end{equation}
Los péndulos se mueven como si el muelle no existiera, y su periodo
$T_0=2\pi/\wz$ da acceso directo al momento de inercia:
\begin{equation}\label{eq:inercia}
  I=\frac{mgL}{\wz^{2}}=\frac{mgL\,T_0^{2}}{4\pi^{2}} .
\end{equation}

\subsubsection*{Oscilación en contrafase}
Si los soltamos con desviaciones opuestas, $\varphi_1(0)=-\varphi_2(0)$,
entonces $A_0=0$ y
\begin{equation}\label{eq:contrafase}
  \varphi_1(t)=-\varphi_2(t)=\varphi_1(0)\cos(\wu t):
\end{equation}
los péndulos se acercan y se alejan a la vez, con periodo $T_1=2\pi/\wu$, más
corto cuanto mayor es $\ell$.

\subsubsection*{Modo pulsante}
Si uno de los péndulos parte del equilibrio y el otro desplazado,
$\varphi_1(0)=0$ y $\varphi_2(0)=\varphi_0$, se excitan los dos modos con la
misma amplitud ($A_0=\varphi_0$, $A_1=-\varphi_0$). Sumando los cosenos de la
\cref{eq:general},
\begin{equation}\label{eq:pulsante}
  \varphi_1(t)=\varphi_0\,\sin(\wm t)\,\sin(\wpl t),
  \qquad
  \varphi_2(t)=\varphi_0\,\cos(\wm t)\,\cos(\wpl t),
\end{equation}
con
\begin{equation}\label{eq:wpm}
  \omega_{\pm}=\frac{\wu\pm\wz}{2}.
\end{equation}
Cada péndulo oscila rápidamente con la frecuencia \emph{portadora} $\wpl$,
muy próxima a $\wz$ y $\wu$, dentro de una \emph{envolvente} lenta de
frecuencia $\wm$. Las dos envolventes, $|\cos\wm t|$ y $|\sin\wm t|$, están
desfasadas un cuarto de ciclo: cuando uno de los péndulos oscila con toda su
amplitud, el otro está casi quieto, y la energía pasa de uno a otro una y
otra vez (\cref{fig:pulsante-teo}). El rozamiento, que no hemos incluido,
amortigua poco a poco todo el proceso. La forma de la señal es la de un
batido, y los apuntes advierten de un detalle que importa para medir:

\begin{referencia}[Frecuencia del batido]{FIV}\label{ref:batido}
Como la estructura de máximos y mínimos se repite dos veces por cada ciclo
de la frecuencia lenta, la frecuencia de batido es la diferencia entre las
frecuencias de las ondas que lo forman.
\end{referencia}

\begin{observacion}[Qué periodo se mide en el modo pulsante]\label{obs:pulsacion}
La amplitud de cada péndulo, $|\cos\wm t|$ o $|\sin\wm t|$, vuelve a su
máximo cada $\pi/\wm$, no cada $2\pi/\wm$, y se anula también cada
$\pi/\wm$. Por eso, si $\Tb$ es el tiempo entre dos máximos consecutivos de
la envolvente o, lo que es lo mismo, entre dos \emph{nodos} consecutivos
(instantes de amplitud nula), que es el \emph{periodo de pulsación}, y $\Tp$
el tiempo entre dos máximos consecutivos de la portadora,
\begin{equation}\label{eq:periodospuls}
  \wm=\frac{\pi}{\Tb},
  \qquad
  \wpl=\frac{2\pi}{\Tp},
  \qquad
  \omega_{\mathrm{b}}=2\wm=\wu-\wz .
\end{equation}
Usar $2\pi/\Tb$ daría un valor de $\wm$ el doble del correcto.
\end{observacion}

\begin{figure}[ht]
\centering
\begin{tikzpicture}
  \pgfmathsetmacro{\wmT}{0.0872}\pgfmathsetmacro{\wpT}{3.2351}
  \pgfmathsetmacro{\Tbt}{pi/\wmT}
  \begin{groupplot}[group style={group size=1 by 2, vertical sep=8pt,
                                  xlabels at=edge bottom, xticklabels at=edge bottom},
                    datos, width=0.95\textwidth, height=3.9cm,
                    xmin=0, xmax=48, ymin=-1.18, ymax=1.72,
                    ytick={-1,0,1}, yticklabels={$-\varphi_0$,$0$,$\varphi_0$},
                    xlabel={$t$ (\si{\second})}, grid=none, clip=false]
  \nextgroupplot[ylabel={$\varphi_2$}]
    \addplot[azulacero!45, line width=0.45pt, domain=0:48, samples=1400]
      {cos(deg(\wmT*x))*cos(deg(\wpT*x))};
    \addplot[azulprof, line width=1.15pt, dash pattern=on 5pt off 2.5pt, domain=0:48, samples=160]
      {abs(cos(deg(\wmT*x)))};
    \addplot[azulprof, line width=1.15pt, dash pattern=on 5pt off 2.5pt, domain=0:48, samples=160]
      {-abs(cos(deg(\wmT*x)))};
    % periodo de pulsación, por encima de la señal
    \draw[grispizarra, line width=0.45pt, |<->|] (axis cs:0,1.38) -- (axis cs:\Tbt,1.38)
      node[midway, above=0pt, text=azulprof, font=\footnotesize] {$\Tb=\pi/\wm$};
  \nextgroupplot[ylabel={$\varphi_1$}]
    \addplot[turquesa!50, line width=0.45pt, domain=0:48, samples=1400]
      {sin(deg(\wmT*x))*sin(deg(\wpT*x))};
    \addplot[turquesa!60!azulprof, line width=1.15pt, dash pattern=on 5pt off 2.5pt, domain=0:48, samples=160]
      {abs(sin(deg(\wmT*x)))};
    \addplot[turquesa!60!azulprof, line width=1.15pt, dash pattern=on 5pt off 2.5pt, domain=0:48, samples=160]
      {-abs(sin(deg(\wmT*x)))};
    % nodos de la envolvente de phi_1
    \addplot[only marks, mark=*, mark size=2.2pt, mark options={fill=white, draw=azulprof, line width=0.7pt}]
      coordinates {(0,0) (\Tbt,0)};
    \node[anchor=south west, text=azulprof, font=\footnotesize, inner sep=1pt] at (axis cs:0.6,0.42) {nodo};
    \node[anchor=south, text=azulprof, font=\footnotesize, inner sep=1pt] at (axis cs:\Tbt,0.42) {nodo};
    % periodo de la portadora, por encima de la señal
    \draw[grispizarra, line width=0.45pt, |<->|]
      (axis cs:{(pi/2+16*pi)/\wpT},1.38) -- (axis cs:{(pi/2+18*pi)/\wpT},1.38)
      node[right=2pt, text=azulprof, font=\footnotesize] {$\Tp=2\pi/\wpl$};
  \end{groupplot}
\end{tikzpicture}
\caption{Modo pulsante según la \cref{eq:pulsante}, con las frecuencias
medidas para $\ell_3=\SI{45}{\centi\meter}$ ($\wz=\SI{3,148}{\radian\per\second}$,
$\wu=\SI{3,322}{\radian\per\second}$) y sin rozamiento. En trazo grueso
discontinuo, las envolventes $\pm|\cos\wm t|$ y $\pm|\sin\wm t|$; en claro,
la portadora. Cuando un péndulo alcanza su máxima amplitud, el otro está en
reposo (nodo de su envolvente). El periodo de pulsación $\Tb$ separa dos
máximos, o dos nodos, consecutivos de la misma envolvente.}
\label{fig:pulsante-teo}
\end{figure}

\subsection{Dependencia con la longitud de acoplamiento}\label{sec:lineales}
La \cref{eq:modos} es \emph{exactamente} lineal en $\ell^{2}$ si la
escribimos para el cuadrado de la frecuencia. Para $\wpl$ y $\wm$ hay que
suponer además que el acoplamiento es débil, $\Omega\ll\wz$, y desarrollar
la raíz a primer orden~\cite{GUION-P3}:
$\wu\simeq\wz\bigl[1+k\ell^{2}/(mgL)\bigr]$. Llevándolo a la \cref{eq:wpm}
quedan las tres relaciones que usaremos:
\begin{equation}\label{eq:lineales}
  \wu^{2}=\wz^{2}+\frac{2k}{I}\,\ell^{2},
  \qquad
  \wpl\simeq\wz+\frac{\wz k}{2mgL}\,\ell^{2},
  \qquad
  \wm\simeq\frac{\wz k}{2mgL}\,\ell^{2}.
\end{equation}

\begin{estimacion}[Tres caminos hacia $k$]\label{est:k}
Las tres magnitudes son rectas $y=a\,\ell^{2}+b$. Conocidos $I$ y $\wz$ por el
modo en fase, cada pendiente da la constante del muelle:
\begin{equation}\label{eq:k}
  k=\frac{I\,a_{1}}{2}\quad\text{(de $\wu^{2}$)},
  \qquad
  k=\frac{2mgL\,a_{\pm}}{\wz}\quad\text{(de $\wpl$ o de $\wm$)},
\end{equation}
y las ordenadas sirven de comprobación: deben valer $\wz^{2}$, $\wz$ y $0$,
respectivamente.
\end{estimacion}

\begin{observacion}[Límites del modelo]\label{obs:limites}
Todo lo anterior supone ángulos pequeños, un muelle ideal (lineal y sin
masa), rozamiento despreciable y dos péndulos idénticos. Además, las
relaciones para $\wpl$ y $\wm$ son aproximadas. Con nuestros datos,
$2k\ell^{2}/(mgL)=\wu^{2}/\wz^{2}-1$ va de \num{0,035} para $\ell_1$ a
\num{0,17} para $\ell_5$, y el término de segundo orden despreciado en $\wm$
es del orden de una cuarta parte de esa cantidad: hasta un \SI{4}{\percent}
para el acoplamiento más fuerte. La relación para $\wu^{2}$, en cambio, no
requiere aproximación alguna.
\end{observacion}

% =========================================================================
\seccion{Procedimiento experimental}{Cómo medimos}\label{sec:procedimiento}
% =========================================================================
Seguimos la metodología del guion~\cite{GUION-P3}, con las herramientas de
análisis de PHYWE Measure~\cite{PHYWE} que describimos a continuación.

\subsection{Registro de las oscilaciones}\label{sec:registro}
Con la fuente, el ordenador y la unidad Cobra-3 encendidos, abrimos PHYWE
Measure y elegimos el modo de medida \emph{Escritor universal}, que registra
la tensión de los dos canales frente al tiempo, sin modificar las opciones de
medida que trae configuradas. Para cada longitud de acoplamiento:
\begin{enumerate}
  \item Situamos las dos pinzas del muelle a la misma altura y medimos $\ell$
        con la cinta métrica desde el eje de giro.
  \item Generamos cada modo con sus condiciones iniciales: los dos péndulos
        desplazados hacia el mismo lado y soltados a la vez (en fase); hacia
        lados opuestos (contrafase); o uno sujeto en la vertical y el otro
        desplazado, soltando ambos a la vez (pulsante). Siempre con
        amplitudes pequeñas y sin dar impulso.
  \item Registramos un número suficiente de oscilaciones: unos
        \SIrange{15}{25}{\second} en los modos en fase y contrafase, y entre
        uno y dos minutos en el pulsante, cuya envolvente es mucho más
        lenta. Guardamos cada medida con un nombre que identifica modo y
        longitud.
\end{enumerate}
Repetimos el procedimiento para $\ell=\SIlist{25;35;45;55;65}{\centi\meter}$.
Como $\wz$ no depende de $\ell$, el modo en fase lo registramos solo para las
dos primeras longitudes, que bastan para determinarlo.

\subsection{Suavizado de la señal}\label{sec:suavizado}
La señal de los sensores no llega limpia: el conversor de la unidad trabaja
en escalones de tensión y hay ruido eléctrico, de modo que cerca de cada
máximo la curva presenta pequeños dientes. Si buscáramos los máximos
directamente, el programa confundiría esos dientes con picos. Por eso
aplicamos antes la función \emph{Suavizar} del menú \emph{Análisis}, que
sustituye cada punto por un promedio de sus vecinos: el ruido desaparece y
la posición de los máximos, que es lo que nos interesa, apenas cambia.
Conviene no exagerar el suavizado, porque un promedio demasiado ancho achata
los picos y puede desplazarlos ligeramente.

\subsection{Análisis de picos}\label{sec:picos}
Con la señal suavizada, \emph{Análisis\,$\to$\,Análisis de pico} localiza los
máximos de uno de los canales (\emph{eje izquierdo} o \emph{derecho}). La
barra de \emph{Tolerancia} decide qué se considera un pico: hacia
\emph{menor} detecta también oscilaciones pequeñas; hacia \emph{mayor} solo
las más pronunciadas. Para cada pico el programa devuelve una tabla con los
instantes de entrada, máximo y salida, la altura en voltios y el área
(\figpanel{fig:picos}{a}), y con \emph{Mostrar resultados} los señala sobre la
gráfica (\figpanel{fig:picos}{b}). De esa tabla tomamos la columna
\emph{máximo} y la altura del pico: son los datos de las
\cref{tab:modo1,tab:modo2,tab:modo3}. La primera fila de la
\figpanel{fig:picos}{a}, por ejemplo, es el máximo $t_1=\SI{2592}{\milli\second}$
del péndulo~1 en fase para $\ell_1$.

\begin{figure}[ht]
\centering
\begin{subfigure}[b]{0.49\textwidth}\centering
  \foto{figuras/analisis-picos.jpg}
  \caption{Ventana \emph{Análisis de picos}.}
\end{subfigure}\hfill
\begin{subfigure}[b]{0.49\textwidth}\centering
  \foto{figuras/picos-contrafase.jpg}
  \caption{Máximos señalados sobre la señal.}
\end{subfigure}
\caption{Análisis de picos. (a)~Tabla de resultados para el péndulo~1 en
fase con $\ell_1$, con los instantes de entrada, máximo y salida de cada pico,
su altura y su área. (b)~Contrafase para $\ell_1$: el programa marca cada
máximo del péndulo~2 con su instante y su tensión
($\SI{1784}{\milli\second}$, $\SI{3672}{\milli\second}$, \dots).}
\label{fig:picos}
\end{figure}

\begin{metodo}[Periodo a partir de los máximos]\label{met:periodo}
Con $N$ máximos consecutivos en los instantes $t_1<t_2<\dots<t_N$ hay $N-1$
periodos completos, de modo que
\begin{equation}\label{eq:periodo}
  T=\frac{t_N-t_1}{N-1}.
\end{equation}
Es la media de los $N-1$ intervalos entre máximos, pero no depende de los
instantes intermedios: solo la cometen los dos extremos, y el error de cada
uno se reparte entre todos los periodos. En los modos en fase, contrafase y
en la portadora tomamos $N=6$. La resolución temporal del registro es de
\SI{4}{\milli\second} (todos los instantes son múltiplos de ese valor), así que
la que corresponde al periodo es unas cinco veces menor.
\end{metodo}

Al revisar los intervalos entre máximos encontramos un único valor anómalo:
el cuarto máximo del péndulo~2 en fase para $\ell_2$, anotado como
\SI{19824}{\milli\second}, rompía la secuencia con saltos de \num{2568} y
\SI{1468}{\milli\second}, frente a los $\approx\SI{2000}{\milli\second}$ del resto.
Se trata de una errata de transcripción (cifras intercambiadas); el valor
correcto, \SI{19284}{\milli\second}, devuelve los intervalos a \num{2028} y
\SI{2008}{\milli\second}. Por la \cref{eq:periodo}, la errata no afectaba al
periodo, pero la corregimos en la \cref{tab:modo1}.

\subsection{El caso pulsante}\label{sec:pulsante}
El modo pulsante merece un comentario aparte, porque en él conviven dos
escalas de tiempo muy distintas y cada una se mide de una manera.

\subsubsection{Cómo se ve en la pantalla}
La \cref{fig:pulsante-fotos} muestra tres registros. Para $\ell_1$ el
acoplamiento es tan débil que en dos minutos apenas cabe una pulsación: la
amplitud del péndulo se apaga hacia la mitad del registro y vuelve a crecer
(\figpanel{fig:pulsante-fotos}{a}). Para longitudes mayores caben varias
(\figpanel{fig:pulsante-fotos}{b}). Con los dos canales a la vista
(\figpanel{fig:pulsante-fotos}{c}) se aprecia el intercambio de energía:
donde la señal azul es más ancha, la roja se estrecha, y viceversa.

\begin{figure}[ht]
\centering
\begin{subfigure}[b]{0.325\textwidth}\centering
  \foto{figuras/pulsante-L1.jpg}
  \caption{$\ell_1$: una pulsación.}
\end{subfigure}\hfill
\begin{subfigure}[b]{0.325\textwidth}\centering
  \foto{figuras/pulsante-varios.jpg}
  \caption{Varias pulsaciones.}
\end{subfigure}\hfill
\begin{subfigure}[b]{0.325\textwidth}\centering
  \foto{figuras/pulsante-dos-canales.jpg}
  \caption{Los dos péndulos.}
\end{subfigure}
\caption{Registros del modo pulsante. (a)~Para $\ell_1=\SI{25}{\centi\meter}$
la envolvente es muy lenta ($\Tb\approx\SI{105}{\second}$) y en el registro
cabe una sola pulsación. (b)~Con un acoplamiento mayor, la envolvente se
repite varias veces en el mismo intervalo de tiempo. (c)~Péndulo~1 (rojo) y
péndulo~2 (azul): la amplitud de uno crece mientras la del otro decrece.}
\label{fig:pulsante-fotos}
\end{figure}

\subsubsection{Dos periodos, dos maneras de medir}
El periodo de la portadora, $\Tp$, se mide exactamente igual que en los
otros dos modos: seis máximos consecutivos y la \cref{eq:periodo}. El periodo
de pulsación, $\Tb$, en cambio, exige localizar la \emph{envolvente}, y la
envolvente no es una curva que el programa registre: solo la vemos a través
de las oscilaciones de la portadora. En lugar de sus máximos, donde la
envolvente es casi plana y varios picos tienen prácticamente la misma altura,
medimos sus \emph{nodos}: los instantes en que la amplitud se anula y la señal
se estrangula sobre la posición de equilibrio, como los nodos de una onda
estacionaria. Si trazamos sobre la gráfica una recta horizontal a la altura
del equilibrio, los nodos son los puntos de la señal que quedan sobre ella
(\cref{fig:picos-envolvente}). Como $|\sin\wm t|$ se anula cada $\pi/\wm$, la
separación entre nodos consecutivos es el mismo periodo de pulsación de la
\cref{obs:pulsacion}, y $\wm=\pi/\Tb$ sigue siendo válida. Para $\ell_1$, por
ejemplo, los dos nodos registrados quedan en
\SIlist{2768;107856}{\milli\second}. Su altura, \SI{-0,22}{\volt}, es la del
equilibrio del canal, muy distinta de la de los máximos de la portadora, en
torno a \SI{-0,13}{\volt} (\cref{tab:modo3}).

\begin{figure}[ht]
\centering
\foto[0.62\textwidth]{figuras/picos-envolvente.jpg}
\caption{Análisis de picos sobre el registro pulsante de $\ell_1$. El programa
marca los picos de la portadora con su instante y su tensión. Los nodos de la
envolvente son los puntos que quedan sobre la línea de equilibrio, a
\SI{-0,22}{\volt}, como el de \SI{2768}{\milli\second} al principio del
registro, donde la amplitud es mínima.}
\label{fig:picos-envolvente}
\end{figure}

Aquí cambia la naturaleza de la incertidumbre, tal como anotamos en el propio
laboratorio: \emph{en $\wm$, la incertidumbre pasa a ser la distancia entre
los máximos}. El nodo de la envolvente cae, en general, entre dos
oscilaciones de la portadora, y la señal solo nos lo deja ver a través de
ellas. La indeterminación en su posición ya no es la del aparato, de unos
pocos milisegundos, sino del orden de la separación entre máximos de la
portadora, $\Tp\approx\SI{2}{\second}$: unas quinientas veces mayor
(\cref{fig:incertidumbre}).

\begin{figure}[ht]
\centering
\begin{tikzpicture}
  \pgfmathsetmacro{\wmT}{0.0886}\pgfmathsetmacro{\wpT}{3.2535}
  \pgfmathsetmacro{\ds}{0.6}
  \begin{axis}[datos, width=0.66\textwidth, height=5.3cm,
               xmin=-9.5, xmax=9.5, ymin=-1.0, ymax=1.3,
               ytick={0}, yticklabels={equilibrio},
               xtick={-8,-4,0,4,8},
               xlabel={$t-t_{\text{nodo}}$ (\si{\second})}, ylabel={$\varphi_1$},
               grid=none, axis lines=left, clip=false]
    % línea de equilibrio
    \draw[grispizarra!60, line width=0.5pt] (axis cs:-9.5,0) -- (axis cs:9.5,0);
    % portadora y envolvente
    \addplot[azulacero!50, line width=0.5pt, domain=-9.5:9.5, samples=700]
      {sin(deg(\wmT*x))*cos(deg(\wpT*(x+\ds)))};
    \addplot[azulprof, line width=1.1pt, dash pattern=on 5pt off 2.5pt, domain=-9.5:9.5, samples=120]
      {abs(sin(deg(\wmT*x)))};
    \addplot[azulprof, line width=1.1pt, dash pattern=on 5pt off 2.5pt, domain=-9.5:9.5, samples=120]
      {-abs(sin(deg(\wmT*x)))};
    % extremos de la portadora: lo único que vemos de la envolvente
    \addplot[puntoA, mark size=1.6pt, domain=-5:5, samples=11]
      ({pi/\wpT*x-\ds},{sin(deg(\wmT*(pi/\wpT*x-\ds)))*cos(deg(pi*x))});
    % nodo
    \addplot[only marks, mark=*, mark size=2.6pt, mark options={fill=white, draw=acento, line width=0.9pt}]
      coordinates {(0,0)};
    \draw[acento, line width=0.4pt] (axis cs:-0.15,0.06) -- (axis cs:-2.6,1.02)
      node[anchor=south east, inner sep=1pt, font=\footnotesize] {nodo de la envolvente};
    % separación entre máximos de la portadora, junto al nodo
    \draw[grispizarra, line width=0.45pt, |<->|]
      (axis cs:{-\ds},0.42) -- (axis cs:{2*pi/\wpT-\ds},0.42)
      node[midway, above=0pt, text=azulprof, font=\footnotesize] {$\Tp$};
  \end{axis}
\end{tikzpicture}
\caption{Incertidumbre en los nodos de la envolvente. Detalle de un nodo con
los parámetros de $\ell_3$: la envolvente (trazo grueso discontinuo) solo se
manifiesta a través de los extremos de la portadora (puntos), separados medio
periodo, y el nodo verdadero puede caer en cualquier punto entre dos
oscilaciones. Su incertidumbre deja de ser la resolución del aparato
(\SI{4}{\milli\second}) y pasa a ser del orden de la separación entre máximos
de la portadora, $\Tp$.}
\label{fig:incertidumbre}
\end{figure}

\subsubsection{Más acoplamiento, más pulsaciones, más precisión}
Con esta limitación, la precisión de $\Tb$ mejora a medida que aumentamos
$\ell$, por dos razones (\cref{fig:precision}).
\begin{itemize}
  \item \emph{Caben más pulsaciones en el registro.} Al crecer $\ell$ la
        envolvente es más rápida: $\Tb$ baja de unos \SI{105}{\second} para
        $\ell_1$ a menos de \SI{18}{\second} para $\ell_5$. Aun con registros
        más cortos, pasamos de dos nodos de la envolvente a cuatro, y la
        indeterminación de los extremos, del orden de $\Tp$, se reparte entre
        más intervalos, como en la \cref{eq:periodo}
        (\figpanel{fig:precision}{a}).
  \item \emph{El nodo está mejor definido.} Con una lectura de tensión de
        resolución $\varepsilon$, los instantes en que la amplitud
        $A\,|\sin\wm t|$ no supera $\varepsilon$ son indistinguibles del
        nodo. Esa ventana dura
        \begin{equation}\label{eq:ventana}
          \Delta t_{\text{nodo}}=\frac{2}{\wm}\arcsin\frac{\varepsilon}{A}
          \simeq\frac{2\varepsilon}{A\,\wm},
        \end{equation}
        inversamente proporcional a $\wm$: pasa de unos \SI{2,8}{\second}
        para $\ell_1$ a medio segundo para $\ell_5$, y desde $\ell_2$ queda
        por debajo de $\Tp$, de modo que la incertidumbre la fija la
        separación entre máximos de la portadora, como decía la nota
        (\figpanel{fig:precision}{b}). En un máximo, donde la envolvente es
        plana, la ventana equivalente,
        $\Delta t_{\text{máx}}=(2/\wm)\arccos(1-\varepsilon/A)$, sería unas
        siete veces mayor: de ahí que midamos nodos y no máximos.
\end{itemize}
Por eso conviene, en general, registrar tantas oscilaciones como sea
posible. Para la envolvente seguimos además un criterio común: cuando el
registro contiene tres o más nodos, descartamos el primero. En el péndulo que
parte del reposo, ese nodo coincide con el instante en que lo soltamos:
mientras lo sujetamos la amplitud ya es nula, y el arranque, nunca perfecto,
contamina su posición. Es visible en $\ell_4$, cuyo primer intervalo, de
\SI{23,5}{\second}, se aparta claramente de los dos siguientes, de
\SI{26,1}{\second}.

\begin{figure}[ht]
\centering
\begin{tikzpicture}
  \begin{groupplot}[group style={group size=2 by 1, horizontal sep=1.6cm},
                    datos, height=5.2cm]
  \nextgroupplot[width=0.55\textwidth, xmin=0, xmax=110, ymin=-2.35, ymax=1.5,
                 xlabel={$t$ (\si{\second})}, ytick=\empty, grid=none,
                 axis y line=none, clip=false,
                 title={\footnotesize\scshape (a)\; registros de \SI{110}{\second}},
                 title style={text=grispizarra}]
    % l1: w- = 0.02989, w+ = 3.1836 ; l5: w- = 0.17651, w+ = 3.2931 (péndulo que parte del reposo)
    \draw[grispizarra!40, line width=0.4pt] (axis cs:0,0.25) -- (axis cs:110,0.25);
    \draw[grispizarra!40, line width=0.4pt] (axis cs:0,-1.45) -- (axis cs:110,-1.45);
    \addplot[azulacero!55, line width=0.4pt, domain=0:110, samples=2200]
      {0.25+0.5*sin(deg(0.02989*x))*sin(deg(3.1836*x))};
    \addplot[azulprof, line width=0.9pt, dash pattern=on 4pt off 2pt, domain=0:110, samples=150]
      {0.25+0.5*abs(sin(deg(0.02989*x)))};
    \addplot[turquesa!55, line width=0.4pt, domain=0:110, samples=2200]
      {-1.45+0.5*sin(deg(0.17651*x))*sin(deg(3.2931*x))};
    \addplot[turquesa!60!azulprof, line width=0.9pt, dash pattern=on 4pt off 2pt, domain=0:110, samples=250]
      {-1.45+0.5*abs(sin(deg(0.17651*x)))};
    \addplot[only marks, mark=*, mark size=1.9pt, mark options={fill=white, draw=azulprof, line width=0.7pt}]
      coordinates {(0,0.25) (105.1,0.25)};
    \addplot[only marks, mark=*, mark size=1.9pt, mark options={fill=white, draw=turquesa!60!azulprof, line width=0.7pt}]
      coordinates {(0,-1.45) (17.8,-1.45) (35.6,-1.45) (53.4,-1.45) (71.2,-1.45) (89.0,-1.45) (106.8,-1.45)};
    \node[anchor=south west, text=azulprof, font=\footnotesize, inner sep=1pt] at (axis cs:0,0.85) {$\ell_1$: 2 nodos};
    \node[anchor=south west, text=turquesa!60!azulprof, font=\footnotesize, inner sep=1pt] at (axis cs:0,-0.85) {$\ell_5$: 7 nodos};
  \nextgroupplot[width=0.41\textwidth, xmin=20, xmax=70, ymin=0.3, ymax=160, ymode=log,
                 xlabel={$\ell$ (\si{\centi\meter})}, ylabel={tiempo (\si{\second})},
                 xtick={25,35,45,55,65}, ytick={0.5,1,2,5,10,20},
                 yticklabels={{0,5},1,2,5,10,20},
                 title={\footnotesize\scshape (b)\; ventana de indistinción},
                 title style={text=grispizarra},
                 legend style={at={(0.98,0.98)}, anchor=north east, font=\scriptsize}]
    % eps = 0.005 V, A = 0.12 V
    \addplot[puntoB] table[x=l, y expr={2*rad(acos(1-0.005/0.12))/\thisrow{wm}}] {
      l   wm
      25  0.02989
      35  0.05601
      45  0.08861
      55  0.12027
      65  0.17651
    };
    \addlegendentry{máximo}
    \addplot[puntoA] table[x=l, y expr={2*rad(asin(0.005/0.12))/\thisrow{wm}}] {
      l   wm
      25  0.02989
      35  0.05601
      45  0.08861
      55  0.12027
      65  0.17651
    };
    \addlegendentry{nodo, ec.~(\ref*{eq:ventana})}
    \addplot[teorica, mark=none] table[x=l, y=Tp] {
      l   Tp
      25  1.9736
      35  1.9560
      45  1.9312
      55  1.9232
      65  1.9080
    };
    \addlegendentry{$\Tp$}
  \end{groupplot}
\end{tikzpicture}
\caption{Por qué gana precisión el modo pulsante al aumentar $\ell$.
(a)~Señal del péndulo que parte del reposo (\cref{eq:pulsante}) para $\ell_1$
y $\ell_5$, con sus frecuencias medidas, durante un mismo intervalo de
\SI{110}{\second}: con el acoplamiento débil cabe una sola pulsación; con el
fuerte, seis. Los círculos son los nodos. (b)~Duración de la ventana en la que
la amplitud no se distingue de la del nodo (\cref{eq:ventana}) o de la del
máximo, con $\varepsilon=\SI{0,005}{\volt}$, media unidad de la lectura del
programa, y $A\approx\SI{0,12}{\volt}$, estimada en la pantalla; en
discontinua, la separación entre máximos de la portadora, $\Tp$. Escala
vertical logarítmica.}
\label{fig:precision}
\end{figure}

\subsection{Tratamiento de incertidumbres}\label{sec:incertidumbres}
Seguimos la \emph{Guide to the expression of uncertainty in
measurement}~\cite{GUM}, tal como se recoge en las notas del
laboratorio~\cite{NUM-ARG}. Para una magnitud $y=f(x_1,\dots,x_N)$ con
entradas no correlacionadas, la incertidumbre típica se propaga como
\begin{equation}\label{eq:propagacion}
  u^{2}(y)=\sum_{k=1}^{N}\left(\frac{\partial f}{\partial x_k}\right)^{\!2}u^{2}(x_k),
\end{equation}
y la incertidumbre ampliada al \SI{95}{\percent} es $\ext{y}=\Tcoef\,u(y)$,
con $\Tcoef$ el factor de cobertura de Student para $\nu$ grados de
libertad~\cite{NUM-ARG}; el resultado se expresa como $y\pm\ext{y}$. Las
incertidumbres se dan con una sola cifra significativa, salvo cuando la
primera es un 1, en cuyo caso se conservan dos, y el valor se redondea a la
misma posición decimal~\cite{ERR-MP}.

\begin{pendientebox}[incertidumbres]
Completar con: (i)~la incertidumbre de los instantes de pico en los modos en
fase y contrafase y en la portadora (resolución de \SI{4}{\milli\second} y
efecto del suavizado); (ii)~la de los nodos de la envolvente, del orden de
$\Tp$ (\cref{sec:pulsante}); (iii)~su propagación a $T$, $\wz$, $\wu$,
$\wm$ y $\wpl$ mediante la \cref{eq:propagacion}; (iv)~la de $\ell$, medida
con la cinta métrica; y (v)~el método de regresión lineal y las
incertidumbres de pendiente y ordenada, como en el
\emph{Método} del informe de Steiner.
\end{pendientebox}

\ornamento

% =========================================================================
\newpage
\seccion{Análisis de datos y resultados}{Lo que dicen las medidas}\label{sec:analisis}
% =========================================================================
\subsection{Datos experimentales}\label{sec:datos}
Las \cref{tab:modo1,tab:modo2,tab:modo3} recogen los instantes de los máximos
obtenidos con el análisis de picos, junto con la altura de los picos y el
periodo que resulta de la \cref{eq:periodo}. Las longitudes de acoplamiento
son $\ell=\SIlist{25;35;45;55;65}{\centi\meter}$, con una incertidumbre
\pend{$u(\ell)$} debida a la cinta métrica. La \cref{tab:frecuencias} reúne
los periodos y las frecuencias angulares de cada longitud, que son los datos
de partida de los ajustes.

% Generado por scripts/tablas.py a partir de Practica3_pendulos_acoplados.xlsx. No editar a mano.

\begin{table}[ht]
\centering
\caption{Modo en fase: instantes de los seis primeros máximos de cada péndulo, altura de los picos y periodo $T=(t_6-t_1)/5$. Este modo solo se registró para $\ell_1$ y $\ell_2$, pues $\omega_0$ no depende de $\ell$.}
\label{tab:modo1}
\renewcommand{\arraystretch}{1.18}\setlength{\tabcolsep}{4.2pt}
\begin{tabular}{@{}l c *{6}{r} c r@{}}
\toprule
$\ell$ & Péndulo & \multicolumn{6}{c}{Tiempo de los máximos, $t_{i}$ (\si{\milli\second})} & $p$ (\si{\volt}) & $T$ (\si{\milli\second}) \\
\cmidrule(lr){3-8}
 & & $t_1$ & $t_2$ & $t_3$ & $t_4$ & $t_5$ & $t_6$ & & \\
\midrule
\multirow{2}{*}{$\ell_{1}=\SI{25}{\centi\meter}$} & 1 & \num{2592} & \num{4588} & \num{6592} & \num{8604} & \num{10644} & \num{12620} & $-$0,13 a $-$0,15 & \num{2005.6} \\
 & 2 & \num{2712} & \num{4720} & \num{6700} & \num{8692} & \num{10676} & \num{12664} & $-$0,06 a $-$0,07 & \num{1990.4} \\
\addlinespace[2pt]
\multirow{2}{*}{$\ell_{2}=\SI{35}{\centi\meter}$} & 1 & \num{1220} & \num{3216} & \num{5204} & \num{7200} & \num{9180} & \num{11176} & $-$0,16 a $-$0,17 & \num{1991.2} \\
 & 2 & \num{13248} & \num{15264} & \num{17256} & \num{19284}\rlap{$^{\dagger}$} & \num{21292} & \num{23232} & $-$0,29 & \num{1996.8} \\
\bottomrule
\end{tabular}
\par\vspace{3pt}{\footnotesize\color{grispizarra}$^{\dagger}$\,En el cuaderno de laboratorio figuraba \SI{19824}{\milli\second}, errata de transcripción (cifras intercambiadas): con ese valor los intervalos vecinos serían de \num{2568} y \SI{1468}{\milli\second}, frente a los $\approx\SI{2000}{\milli\second}$ del resto.\par}
\end{table}

% Generado por scripts/tablas.py a partir de Practica3_pendulos_acoplados.xlsx. No editar a mano.

\begin{table}[ht]
\centering
\caption{Modo en contrafase: instantes de los seis primeros máximos de cada péndulo, altura de los picos y periodo $T=(t_6-t_1)/5$ para las cinco longitudes de acoplamiento.}
\label{tab:modo2}
\renewcommand{\arraystretch}{1.18}\setlength{\tabcolsep}{4.2pt}
\begin{tabular}{@{}l c *{6}{r} c r@{}}
\toprule
$\ell$ & Péndulo & \multicolumn{6}{c}{Tiempo de los máximos, $t_{i}$ (\si{\milli\second})} & $p$ (\si{\volt}) & $T$ (\si{\milli\second}) \\
\cmidrule(lr){3-8}
 & & $t_1$ & $t_2$ & $t_3$ & $t_4$ & $t_5$ & $t_6$ & & \\
\midrule
\multirow{2}{*}{$\ell_{1}=\SI{25}{\centi\meter}$} & 1 & \num{2752} & \num{4708} & \num{6676} & \num{8648} & \num{10616} & \num{12576} & $-$0,15 a $-$0,16 & \num{1964.8} \\
 & 2 & \num{1784} & \num{3672} & \num{5608} & \num{7616} & \num{9556} & \num{11580} & $-$0,30 & \num{1959.2} \\
\addlinespace[2pt]
\multirow{2}{*}{$\ell_{2}=\SI{35}{\centi\meter}$} & 1 & \num{2376} & \num{4320} & \num{6260} & \num{8200} & \num{10132} & \num{12076} & $-$0,13 a $-$0,14 & \num{1940.0} \\
 & 2 & \num{1456} & \num{3376} & \num{5256} & \num{7200} & \num{9108} & \num{11044} & $-$0,28 & \num{1917.6} \\
\addlinespace[2pt]
\multirow{2}{*}{$\ell_{3}=\SI{45}{\centi\meter}$} & 1 & \num{680} & \num{2592} & \num{4500} & \num{6404} & \num{8316} & \num{10212} & $-$0,17 a $-$0,20 & \num{1906.4} \\
 & 2 & \num{1688} & \num{3484} & \num{5404} & \num{7248} & \num{9172} & \num{11068} & $-$0,30 a $-$0,31 & \num{1876.0} \\
\addlinespace[2pt]
\multirow{2}{*}{$\ell_{4}=\SI{55}{\centi\meter}$} & 1 & \num{19516} & \num{21400} & \num{23292} & \num{25168} & \num{27068} & \num{28936} & $-$0,16 & \num{1884.0} \\
 & 2 & \num{9424} & \num{11268} & \num{13168} & \num{15064} & \num{16924} & \num{18824} & $-$0,07 a $-$0,09 & \num{1880.0} \\
\addlinespace[2pt]
\multirow{2}{*}{$\ell_{5}=\SI{65}{\centi\meter}$} & 1 & \num{15132} & \num{16996} & \num{18856} & \num{20720} & \num{22556} & \num{24388} & $-$0,20 a $-$0,21 & \num{1851.2} \\
 & 2 & \num{5272} & \num{7148} & \num{8996} & \num{10812} & \num{12668} & \num{14468} & $-$0,11 & \num{1839.2} \\
\bottomrule
\end{tabular}
\end{table}

% Generado por scripts/tablas.py a partir de Practica3_pendulos_acoplados.xlsx. No editar a mano.

\begin{table}[ht]
\centering
\caption{Modo pulsante. Arriba, instantes de los nodos de la envolvente (amplitud nula, sobre la línea de equilibrio) y periodo de pulsación $T_{\mathrm{b}}=(t_{\text{últ}}-t_{\text{pri}})/n$, con $n$ el número de intervalos empleados (entre paréntesis, el nodo descartado). Abajo, seis máximos consecutivos de la portadora y su periodo $T_{+}=(t_6-t_1)/5$.}
\label{tab:modo3}
\renewcommand{\arraystretch}{1.18}\setlength{\tabcolsep}{4.6pt}
\begin{tabular}{@{}l *{4}{r} c c r@{}}
\toprule
\multicolumn{8}{@{}l}{\scshape\color{azulprof}Envolvente (nodos)} \\
$\ell$ & \multicolumn{4}{c}{$t_{i}$ (\si{\milli\second})} & $p$ (\si{\volt}) & $n$ & $T_{\mathrm{b}}$ (\si{\milli\second}) \\
\cmidrule(lr){2-5}
 & $t_1$ & $t_2$ & $t_3$ & $t_4$ & & & \\
\midrule
$\ell_{1}=\SI{25}{\centi\meter}$ & \num{2768} & \num{107856} & — & — & $-$0,22 & 1 & \num{105088} \\
$\ell_{2}=\SI{35}{\centi\meter}$ & \num{468} & \num{56560} & — & — & $-$0,23 a $-$0,24 & 1 & \num{56092} \\
$\ell_{3}=\SI{45}{\centi\meter}$ & \textcolor{grispizarra!70}{(\num{2248})} & \num{38152} & \num{73608} & — & $-$0,21 a $-$0,22 & 1 & \num{35456} \\
$\ell_{4}=\SI{55}{\centi\meter}$ & \textcolor{grispizarra!70}{(\num{2300})} & \num{25756} & \num{51860} & \num{78000} & $-$0,21 a $-$0,22 & 2 & \num{26122} \\
$\ell_{5}=\SI{65}{\centi\meter}$ & \textcolor{grispizarra!70}{(\num{2256})} & \num{19852} & \num{37516} & \num{55448} & $-$0,22 & 2 & \num{17798} \\
\bottomrule
\end{tabular}
\par\vspace{8pt}
\begin{tabular}{@{}l *{6}{r} c r@{}}
\toprule
\multicolumn{9}{@{}l}{\scshape\color{azulprof}Portadora (máximos consecutivos)} \\
$\ell$ & \multicolumn{6}{c}{$t_{i}$ (\si{\milli\second})} & $p$ (\si{\volt}) & $T_{+}$ (\si{\milli\second}) \\
\cmidrule(lr){2-7}
 & $t_1$ & $t_2$ & $t_3$ & $t_4$ & $t_5$ & $t_6$ & & \\
\midrule
$\ell_{1}=\SI{25}{\centi\meter}$ & \num{40796} & \num{42744} & \num{44708} & \num{46684} & \num{48668} & \num{50664} & $-$0,13 & \num{1973.6} \\
$\ell_{2}=\SI{35}{\centi\meter}$ & \num{18244} & \num{20164} & \num{22172} & \num{24116} & \num{26064} & \num{28024} & $-$0,14 & \num{1956.0} \\
$\ell_{3}=\SI{45}{\centi\meter}$ & \num{14264} & \num{16124} & \num{18060} & \num{20020} & \num{22000} & \num{23920} & $-$0,11 a $-$0,12 & \num{1931.2} \\
$\ell_{4}=\SI{55}{\centi\meter}$ & \num{10236} & \num{12136} & \num{14068} & \num{15968} & \num{17932} & \num{19852} & $-$0,12 a $-$0,14 & \num{1923.2} \\
$\ell_{5}=\SI{65}{\centi\meter}$ & \num{6216} & \num{8108} & \num{10020} & \num{11892} & \num{13856} & \num{15756} & $-$0,15 a $-$0,18 & \num{1908.0} \\
\bottomrule
\end{tabular}
\end{table}

% Generado por scripts/tablas.py a partir de Practica3_pendulos_acoplados.xlsx. No editar a mano.

\begin{table}[ht]
\centering
\caption{Periodos y frecuencias angulares para cada longitud de acoplamiento. $T_1$ es la media de los dos péndulos en contrafase; $\omega_1=2\pi/T_1$, $\omega_{-}=\pi/T_{\mathrm{b}}$ y $\omega_{+}=2\pi/T_{+}$. Las incertidumbres se tratan en la \cref{sec:ajustes}.}
\label{tab:frecuencias}
\renewcommand{\arraystretch}{1.22}\setlength{\tabcolsep}{5pt}
\begin{tabular}{@{}l c c c c c c c c@{}}
\toprule
$\ell$ & $\ell^{2}$ (\si{\meter\squared}) & $T_1$ (\si{\second}) & $\omega_1$ (\si{\radian\per\second}) & $\omega_1^{2}$ (\si{\radian\squared\per\second\squared}) & $T_{\mathrm{b}}$ (\si{\second}) & $\omega_{-}$ (\si{\radian\per\second}) & $T_{+}$ (\si{\second}) & $\omega_{+}$ (\si{\radian\per\second}) \\
\midrule
$\ell_{1}$ & \num{0.0625} & \num{1.9620} & \num{3.2024} & \num{10.256} & \num{105.09} & \num{0.0299} & \num{1.9736} & \num{3.1836} \\
$\ell_{2}$ & \num{0.1225} & \num{1.9288} & \num{3.2576} & \num{10.612} & \num{56.09} & \num{0.0560} & \num{1.9560} & \num{3.2123} \\
$\ell_{3}$ & \num{0.2025} & \num{1.8912} & \num{3.3223} & \num{11.038} & \num{35.46} & \num{0.0886} & \num{1.9312} & \num{3.2535} \\
$\ell_{4}$ & \num{0.3025} & \num{1.8820} & \num{3.3386} & \num{11.146} & \num{26.12} & \num{0.1203} & \num{1.9232} & \num{3.2670} \\
$\ell_{5}$ & \num{0.4225} & \num{1.8452} & \num{3.4052} & \num{11.595} & \num{17.80} & \num{0.1765} & \num{1.9080} & \num{3.2931} \\
\bottomrule
\end{tabular}
\end{table}


\subsection{Frecuencia natural y momento de inercia}\label{sec:inercia}
El modo en fase nos da cuatro medidas del periodo $T_0$: una por péndulo y
longitud (\cref{tab:omega0}). Las cuatro coinciden dentro de un
\SI{0,8}{\percent}, como debe ser si $\wz$ es propia de cada péndulo y no del
acoplamiento, y su media es
\begin{equation}\label{eq:T0}
  \mean{T}_0=\valTcero\ \si{\second}
  \qquad\Longrightarrow\qquad
  \wz=\frac{2\pi}{\mean{T}_0}=\valWcero\ \si{\radian\per\second},
  \qquad
  \wz^{2}=\valWceroDos\ \si{\radian\squared\per\second\squared}.
\end{equation}

% Generado por scripts/tablas.py a partir de Practica3_pendulos_acoplados.xlsx. No editar a mano.

\begin{table}[ht]
\centering
\caption{Periodos del modo en fase y frecuencia propia de cada medida. La última fila recoge el periodo medio y la frecuencia $\omega_0=2\pi/\mean{T}_0$ que empleamos en adelante.}
\label{tab:omega0}
\renewcommand{\arraystretch}{1.22}
\begin{tabular}{@{}l c c c@{}}
\toprule
Longitud & Péndulo & $T_0$ (\si{\second}) & $2\pi/T_0$ (\si{\radian\per\second}) \\
\midrule
$\ell_1=\SI{25}{\centi\meter}$ & 1 & \num{2.0056} & \num{3.1328} \\
$\ell_1=\SI{25}{\centi\meter}$ & 2 & \num{1.9904} & \num{3.1567} \\
$\ell_2=\SI{35}{\centi\meter}$ & 1 & \num{1.9912} & \num{3.1555} \\
$\ell_2=\SI{35}{\centi\meter}$ & 2 & \num{1.9968} & \num{3.1466} \\
\midrule
\multicolumn{2}{@{}l}{\scshape Media} & \num{1.9960} & \num{3.1479} \\
\bottomrule
\end{tabular}
\end{table}


\begin{calculo}[Momento de inercia]\label{calc:I}
Con $m=\SI{1}{\kilogram}$, $L=\SI{1}{\meter}$ y $g=\SI{9,81}{\meter\per\second\squared}$,
la \cref{eq:inercia} da
\begin{equation}\label{eq:Ivalor}
  I=\frac{mgL}{\wz^{2}}
   =\frac{\SI{1}{\kilogram}\cdot\SI{9,81}{\meter\per\second\squared}\cdot\SI{1}{\meter}}
         {\valWceroDos\ \si{\radian\squared\per\second\squared}}
   =\valI\ \si{\kilogram\meter\squared}.
\end{equation}
\end{calculo}

\begin{destacado}[Cuestión 1]
$\wz=\valWcero\ \si{\radian\per\second}$
\quad\textperiodcentered\quad
$I=\valI\ \si{\kilogram\meter\squared}$
\;\pend{$\pm\,U(I)$, al \SI{95}{\percent}}
\end{destacado}

\begin{observacion}[Comprobaciones]\label{obs:comprobaciones}
Hemos contrastado el resultado por dos vías independientes de la hoja de
cálculo. Promediar las cuatro frecuencias $2\pi/T_0$ en lugar de los
periodos da \SI{3,1479}{\radian\per\second}, idéntico hasta la cuarta cifra
decimal. Y si en vez de la \cref{eq:periodo} obtenemos cada periodo como la
pendiente de la recta $t_i=t_1+(i-1)\,T$ ajustada a los seis máximos, la media
pasa a ser \SI{1,9974}{\second} e $I=\SI{0,9914}{\kilogram\meter\squared}$: una
diferencia del \SI{0,14}{\percent}, muy por debajo de la dispersión entre las
cuatro medidas.
\end{observacion}

\begin{observacion}[Un péndulo casi simple]\label{obs:simple}
El valor obtenido es muy próximo al de una masa puntual de \SI{1}{\kilogram}
a \SI{1}{\meter} del eje, $mL^{2}=\SI{1}{\kilogram\meter\squared}$. Dicho de otro
modo, el periodo medido, $\mean{T}_0=\SI{1,996}{\second}$, es casi el de un
péndulo simple de un metro, $2\pi\sqrt{L/g}=\SI{2,006}{\second}$: la masa de la
varilla y las dimensiones de la masa lenticular apenas alteran el comportamiento, y
la pequeña diferencia sugiere que el centro de masas queda algo por encima
de $L=\SI{1}{\meter}$.
\end{observacion}

\begin{pendientebox}[cuestión 1]
Incertidumbre de $\mean{T}_0$ (tipo~A con $s/\sqrt{4}$ y $\nu=3$, y tipo~B por la
resolución) y su propagación a $\wz$ e $I$ mediante la
\cref{eq:propagacion}: $u(\wz)=2\pi\,u(\mean{T}_0)/\mean{T}_0^{2}$ y
$u(I)/I=2\,u(\wz)/\wz$. La hoja de cálculo da, como estimación preliminar,
$s/\sqrt{4}=\SI{0,0035}{\second}$.
\end{pendientebox}

\subsection{Ajustes lineales y constante del muelle}\label{sec:ajustes}
\begin{pendientebox}[cuestión 2]
Representar $\wu^{2}$, $\wm$ y $\wpl$ frente a $\ell^{2}$ con los datos de la
\cref{tab:frecuencias} y sus barras de incertidumbre; ajustar las tres rectas
$y=a\,\ell^{2}+b$; dar pendientes y ordenadas con sus incertidumbres, y
deducir $k$ con la \cref{eq:k}. Comparar cada ordenada con su valor
teórico ($\wz^{2}=\valWceroDos\ \si{\radian\squared\per\second\squared}$,
$0$ y $\wz=\valWcero\ \si{\radian\per\second}$).
\end{pendientebox}

\begin{figure}[ht]
\centering
\huecofigura[4.6cm]{Gráfica: $\wu^{2}$ frente a $\ell^{2}$ con su recta de ajuste.}\\[6pt]
\huecofigura[4.6cm]{Gráfica: $\wm$ y $\wpl$ frente a $\ell^{2}$ con sus rectas de ajuste.}
\caption{\pendiente{Pie de figura de los ajustes lineales.}}
\label{fig:ajustes}
\end{figure}

\begin{table}[ht]
\centering
\caption{Parámetros de los ajustes $y=a\,\ell^{2}+b$ y constante del muelle
que se deduce de cada uno.}
\label{tab:ajustes}
\renewcommand{\arraystretch}{1.3}
\begin{tabular}{@{}l c c c c c@{}}
\toprule
Ajuste & $a\pm U(a)$ & $b\pm U(b)$ & Ordenada teórica & $r$ & $k\pm U(k)$ (\si{\newton\per\meter}) \\
\midrule
$\wu^{2}$ frente a $\ell^{2}$ & \pend{} & \pend{} & $\wz^{2}$ & \pend{} & \pend{} \\
$\wm$ frente a $\ell^{2}$     & \pend{} & \pend{} & $0$       & \pend{} & \pend{} \\
$\wpl$ frente a $\ell^{2}$    & \pend{} & \pend{} & $\wz$     & \pend{} & \pend{} \\
\bottomrule
\end{tabular}
\end{table}

\subsection{Comparación de los tres métodos}\label{sec:comparacion}
\begin{pendientebox}[cuestión 3]
Comparar los tres valores de $k$ (compatibilidad de sus intervalos al
\SI{95}{\percent}) y discutir cuál es más preciso y por qué. Puntos que
pueden ayudar: la relación de $\wu^{2}$ es exacta, mientras que las de
$\wm$ y $\wpl$ son aproximadas (\cref{obs:limites}); $\wpl$ cambia muy poco
con $\ell$, de modo que su pendiente es pequeña frente a la dispersión; y
$\wm$ procede de la envolvente, cuya incertidumbre es del orden de $\Tp$
(\cref{sec:pulsante}).
\end{pendientebox}

\ornamento

% =========================================================================
\newpage
\seccion{Discusión y conclusiones}{Lo que aprendemos}\label{sec:discusion}
% =========================================================================
\begin{pendientebox}[conclusiones]
Redactar las conclusiones sobre $\wz$ e $I$ (\cref{sec:inercia}), los ajustes y
los valores de $k$ (\cref{sec:ajustes}), la comparación entre métodos
(\cref{sec:comparacion}) y las fuentes de error: amplitudes finitas,
rozamiento, masa del muelle, condición inicial del modo pulsante y número de
pulsaciones registradas.
\end{pendientebox}

% =========================================================================
\newpage
%                              BIBLIOGRAFÍA
% =========================================================================
\nocite{TIPLER,HUGHES,FIV-T4}
\printbibliography[title={Referencias}]

\end{document}
